%% September 03, 2004
%% Version 1.0
\documentclass[12pt]{report}
\usepackage{amsmath,latexsym,epic,eepic,pstcol,comment}%,framed}
\usepackage[dvips,hyperref]{hyperref}

\usepackage{makeidx}
\makeindex

%%\renewcommand{\thefigure}{\Roman{chapter}.\arabic{figure}}
\renewcommand{\thefigure}{\arabic{chapter}.\arabic{figure}}

\newcommand{\code}[1]{\texttt{#1}}
\newcommand{\filename}[1]{\texttt{#1}}
%% \newcommand{\code}[1]{\texttt{#1}}

%% \setlength{\topmargin}{0.5in}
\newcommand{\ePiX}{\code{ePiX}}
\newcommand{\ext}{\code{ePiX\_ext}}
\newcommand{\epix}{\code{epix}}
\newcommand{\eepic}{\code{eepic}}
\newcommand{\C}{\code{C}}
\newcommand{\CXX}{\code{C++}}
\newcommand{\pyepix}{\code{Pyepix}}

\newcommand{\bi}{\begin{itemize}}
\newcommand{\ei}{\end{itemize}}

\title{ \ePiX\ Tutorial and  \\ Reference Manual}
\author{Andrew D. Hwang \\ 
  Department\ of Math and CS \\ 
  College of the Holy Cross}
\date{Version 1.0,  September, 2004}

\begin{document}

\maketitle

\tableofcontents

\chapter{Introduction}

\ePiX, a collection of batch-oriented utilities for *nix, creates
mathematically accurate line figures, plots, and movies using
easy-to-learn syntax. \LaTeX~and \code{dvips} comprise the
typographical rendering engine, while \code{ImageMagick} is used to
create bitmapped images and animations.  The user interface resembles
that of \LaTeX\ itself: You prepare a short scene description in a
text editor, then ``compile'' the input file into a picture. Default
output formats are \eepic\ (a plain text enhancement to the \LaTeX\
picture environment), \code{eps}, \code{pdf}, \code{png},
and~\code{mng}. 

\ePiX's strengths include:
 
\begin{itemize}
\item Quality of output: \ePiX\ creates mathematically accurate,
  publication-quality figures whose appearance matches that of \LaTeX.
  Typography may be put in a figure as easily as in an ordinary
  \LaTeX\ document.
  
\item Ease of use: Figure objects and their attributes are specified
  by simple, descriptive commands.
  
\item Flexibility: In \ePiX, an object is described by its attributes
  and Cartesian location; as in \LaTeX, printed appearance is
  determined when the figure is compiled. A well-designed figure can
  be altered dramatically, yet precisely, with minor changes to the
  input file.
  
\item Power and extendability: \ePiX\ inherits the power of \CXX\ as a
  programming language; variables, data structures, loops, and
  recursion can be used to draw complicated plots and figures with
  just a few lines of input. External code can be incorporated in an
  \ePiX\ figure with a command line option or by using a Makefile.
  
\item Economy of storage and transmission: For a document containing
  many figures, a compressed tar file of the \LaTeX\ sources and
  \ePiX\ files is typically a few percent the size of the compressed
  PostScript file.
  
\item License: \ePiX\ is \emph{Free Software}.\index{Free software}
  You are granted the right to use the program for whatever purpose,
  and to inspect, modify, and re-distribute the source code, so long
  as you do not restrict the rights of others to do the same. In
  short, the license is similar to the terms under which theorems are
  published.
  
\end{itemize}


The relationship of \ePiX\ to a graphical drawing program is analogous
to the relationship between \LaTeX\ and a word processor; \ePiX\ 
facilitates logical structuring of mathematical figures. Though \ePiX\ 
makes a few stylistic defaults to streamline the creation of simple
figures, it imposes no internal restrictions on the contents or
appearance of a figure; aesthetic and practical decisions are left to
the user.

This manual is meant to be read in stages rather than ``cover to
cover''. If you are a:

\bi

\item Potential user, you may wish to skip immediately to ``Software
Dependencies'' before investing additional time.

\item New user, please proceed from here until you have enough
  understanding to play with the software, then experiment with the
  samples files while reading Chapter~\ref{chapter:started}, or return
  to the manual as needed.

\item More advanced user, browse at will, probably starting with
  Chapter~\ref{chapter:ref-man}.

\ei

In any case, please don't hesitate to contact the author or join the
mailing list if you have questions or comments (good or bad) about the
software or manual, or if you are willing and able to join in
developmment.

Under the philosophy that people learn most easily when ideas are
introduced in context, this manual is relatively conversational, and
occasionally redundant (especially between portions meant for readers
at different levels of familiarity). Throughout, you are assumed to be
familiar with \LaTeX\ and basic linear algebra (the description of
points, vectors, lines, and planes in three-dimensional space). Other
material, such as \CXX\ syntax, is introduced as needed.


\section{Software Dependencies}

\index{Installation|(} If you run GNU/Linux, a BSD, or Solaris, you
almost surely have (or can install) all the external software needed
to use \ePiX.  For Mac~OS~X,\index{Mac OS X} you will need the Apple
developer tools, and will want to install an X~server and
the~\code{fink} package manager to build a featureful *nix
environment.  Users of other operating systems, most notably Windows,
face a challenge in running \ePiX, though not an insurmountable one.
The author's (second-hand) Windows-specific knowledge is summarized
below.

``Under the hood'', an input file is successively converted
to~\code{eepic}, \code{dvi}, PostScript, \code{pdf}~or \code{eps}, and
if desired, \code{png}~or \code{mng}.  \ePiX\ comprises a compiled
library written in~\CXX, a \CXX\ header file, and four shell
scripts---\epix, \code{laps}, \code{elaps}, and~\code{flix}---that
automate the various file format conversions.  Each script is written
in GNU~\code{bash}.  Consequently, there are two absolute necessities:
a \CXX\ compiler (preferably~\code{g++}) and \code{bash}. The
script~\code{epix} uses only these programs.  \ePiX\ is primarily a
pre-processor for \LaTeX, but does not absolutely require \LaTeX\ for
normal use. However, without \LaTeX~and Ghostscript
(particularly~\code{dvips}), you can't view or print \code{epix}'s
output files, nor can you run \code{flix} or~\code{elaps}.

A text editor such as \code{emacs} or \code{vim} that facilitates
formatting \C~code is extremely useful for writing input files. \ePiX\
comes with an \code{emacs} mode, written by Jay Belanger, that allows
you to write, compile, and view \ePiX\ figures without leaving
\code{emacs}. 

In the presence of \LaTeX~and Ghostscript, only a few standard
utilities are needed to run \ePiX's ``conversion to
\code{eps}/\code{pdf}'' script~\code{elaps}, namely \code{grep},
\code{sed}, \code{epstopdf}, and \code{ps2epsi}.

Finally, \code{flix} uses ImageMagick's \code{convert} utility to
create \code{png}~images, and to assemble \code{png}s into animated
\code{mng} files. The programs \code{animate}~and \code{display} are
useful for viewing \code{flix}~output.\index{Animation}

Aside from their reliance on specific programs, \ePiX's shell scripts
are written using Unix-style pathnames. Thus, the most straightforward
way to use \ePiX\ is to install a Unix-like environment.


\subsection*{Alternatives for Windows}

\index{Windows|(}
Version~0.8.9 of \ePiX\ has been implemented in Python~2.2 by Andrew
Sterian, making \ePiX\ available on any platform that supports Python,
and without requiring a \CXX~compiler or~\code{bash}.  Python is a
GPL-ed scripting language, and is available with a Windows installer
and detailed instructions.  The easiest alternative for Windows users
is probably to install Python~2.2 or later (if necessary) and
\code{Pyepix}. The \pyepix\ project home page is:
\begin{verbatim}
  http://claymore.engineer.gvsu.edu/~steriana/Python/index.html
\end{verbatim}

Cygwin, and theoretically the de~Lorie tools, can be used to run
\ePiX\ under Windows. The following Cygwin packages are probably
necessary and sufficient: 
\begin{verbatim}
  [] bash       [] GhostScript [] sh-utils
  [] binutils   [] ImageMagick [] teTeX
  [] fileutils  [] make        [] textutils
  [] gcc        [] sed
\end{verbatim}
\begin{comment}
  \filename{bash}, \filename{binutils}, \filename{fileutils},
  \filename{gcc}, \filename{GhostScript}, \filename{ImageMagick},
  \filename{make}, \filename{sed}, \filename{sh-utils},
  \filename{teTeX}, \filename{textutils}.
\end{comment}
The \filename{emacs}~and \filename{gv} packages are highly desirable,
but not absolutely necessary. The author has received scattered
reports of success with Cygwin, but is not sufficiently knowledgeable
to provide substantial support.
\index{Windows|)}


\section{Installation}

\ePiX\ is distributed over the World-Wide Web as source code. Packages
(stable and development) can be found at
\begin{verbatim}
  http://math.holycross.edu/~ahwang/current/ePiX.html
\end{verbatim}
In a web browser, \texttt{shift}-click on a link to download.  The
latest stable release is also on the CTAN mirrors, in the
\code{graphics} directory. There are instructions for downloading an
entire directory; it is not recommended that you download the files
individually.  (Some users of Red Hat have reported file permission
problems when unpacking the CTAN tarballs. If you encounter this
difficulty, please try downloading the sources from the project main
page.)  Unpack the compressed tar file with the appropriate command:
\begin{verbatim}
  tar -zxvf epix-x.y.z_complete.tar.gz
  tar -jxvf epix-x.y.z_complete.tar.bz2
\end{verbatim}
(\code{x.y.z}~is the version number) or, if your \code{tar}
doesn't do decompression,
\begin{verbatim}
  gunzip -c epix-x.y.z_complete.tar.gz | tar -xvf -
  bzcat epix-x.y.z_complete.tar.bz2 | tar -xvf -
\end{verbatim}
\code{cd} to the source directory, which is named \filename{epix-x.y.z}.
The \filename{INSTALL} file contains detailed installation instructions.
If you're impatient, the short of it is \code{make [test]; make install}.
Respectively, these steps build the program, optionally run a test
compile on the included sample files, and install the library, header
file, and shell scripts.

There is an optional package, contributed by Svend Daug{\aa}rd
Pedersen, that supplies extensions for enhanced Cartesian and
logarithmic coordinate systems, and for hatching polygons and planar
regions. To build this package, do \code{make contrib} before
\code{make install}.

In order to use the \filename{contrib} package, an input file must
contain the line ``\code{using namespace ePiX\_contrib;}'' For
documentation, please see the directory
\filename{\$INSTALL/share/epix/tutorial/contrib}.

By default, \ePiX\ installs in subdirectories of \filename{/usr/local};
if you want to install elsewhere, see the \filename{INSTALL} file for
detailed instructions. You may also want to consult
\filename{POST-INSTALL} for information on setting your \code{PATH}
variable so your shell can find \ePiX.  \index{Installation|)}

To re-iterate, \ePiX\ is not a stand-alone program, but consists of a
\C/\CXX\ library, header, and a shell script, and therefore requires a
compiler \emph{for normal use}. The GNU compiler (\code{g++}) and
\CXX\ library are strongly preferred, both because they are used to
develop \ePiX, and because they implement many mathematical features
not specified by \code{ANSI~C}.  \ePiX\ is also reliant on the GNU
shell \code{bash}. Most Unices (and all major GNU/Linux distributions)
have \code{bash} in \filename{/bin/bash}. The \code{INSTALL} file
explains how to cope with \code{bash} that is elsewhere. If you port
\ePiX\ to another shell or operating environment, or package \ePiX\ 
for other systems, please notify the author so that your work can be
linked from the project page and mentioned in the documentation.

Should it come to this, the command \code{make uninstall} will remove
installed components of \ePiX\ from your system. You must be in the
source directory, and may need to log in as \code{root}.


\subsection*{Other Sources}

An RPM \filename{spec} file is maintained by Guido Gonzato. If you run
a GNU/Linux distribution that uses RPM~4.x to manage installed
packages (e.g., a recent version of Red Hat or Mandrake), you can
build and install \ePiX\ with the command (as \code{root})
\begin{verbatim}
rpm -ta epix_complete.tar.gz
\end{verbatim}
An \filename{rpm} source file is available from the project page at
Holy Cross. Separately maintained packages exist for FreeBSD and
Gentoo:
\begin{verbatim}
  http://www.freshports.org/graphics/epix/
  http://packages.gentoo.org/search/?sstring=epix
\end{verbatim}
The very latest source code for Versions 1.x~and 2.x (a.k.a.\ The
Next Generation) is available by CVS from
\begin{verbatim}
  http://savannah.nongnu.org/cgi-bin/viewcvs/epix/epix/
\end{verbatim}
The main project page may migrate to
\begin{verbatim}
  http://savannah.nongnu.org/projects/epix
\end{verbatim}


\subsection*{Development}

There is a (currently very low-traffic) mailing list for all things
\ePiX-related:
\begin{verbatim}
  ahwang-epix@mathcs.holycross.edu
\end{verbatim}
You must be a subscriber in order to post, and a post must have a
non-empty \code{Subject:} line.  To subscribe, send an empty message
to
\begin{verbatim}
  ahwang-epix-subscribe@mathcs.holycross.edu
\end{verbatim}
Please contact the author if you are interested in development.


\chapter{Getting Started}
\label{chapter:started}

This chapter describes the basics of creating figures in \ePiX, and is
written for readers who are familiar with \LaTeX\ but completely new
to~\CXX. No detailed knowledge of~\CXX\ is needed to use \ePiX, only a
bit of grammar (roughly \code{C}~grammar, without pointers) that is
easily absorbed by example.  As you read, please study the sample
files that are distributed with \ePiX. If you have installed in the
default location, the samples are under
\filename{/usr/local/share/epix}. The quickest way to learn is to copy
sample files to a convenient location and experiment by modifying them
and seeing how your changes affect the figure's appearance.

\section{Running \ePiX}
\label{section:running}

The sample and tutorial figures carry the preferred file extension,
\filename{.xp}. To convert \filename{file.xp} to a viewable figure, do
``\filename{elaps file.xp}'' to produce an \code{eps} file, or
``\filename{elaps --pdf file.xp}'' for PDF.  Like \LaTeX, \ePiX\ is
non-interactive, and run from the command line.  Each shell script
accepts a \code{--help} (or~\code{-h}) option that prints a short
summary of available options.  Under~\code{X}, a graphical environment
may be simulated by using \code{emacs} (with Jay Belanger's \ePiX\ 
mode) to edit and compile files, and a previewer such as~\code{gv} to
view figures as they are processed. In~\code{gv}, select ``Watch
file'' from the ``State'' menu to have images update
automatically.\index{Graphical interface}

While \code{elaps} produces output files suitable for immediate
previewing, \code{epix} is far preferable for documents processed with
\LaTeX\ (rather than PDF\LaTeX). The output of~\code{epix} is
\filename{eepic}, an enhancement to the \LaTeX\ picture environment
that allows lines of arbitrary length, slope, and width. An
\filename{eepic} file is typically a few percent the size of a
comparable \filename{eps} file, and is a text file that can be read
(and edited) with only knowledge of \LaTeX.

An \filename{eepic} file is included directly into a \LaTeX\ document,
say \filename{sample.tex}, that contains the line
\begin{verbatim}
\usepackage{epic,eepic,pstcol}
\end{verbatim}
Somewhere in \filename{sample.tex} is the figure itself:
\begin{verbatim}
\begin{figure}[hbt]
  \begin{center}
    \input{example.eepic}
    \caption{A compiled \ePiX\ figure.}
    \label{fig:example}
  \end{center}
\end{figure}
\end{verbatim}
In a text editor, create an \ePiX\ input file for the figure, say
\filename{file.xp}, then issue the command
\begin{verbatim}
epix file.xp example.eepic
\end{verbatim}
File extensions may be omitted, as can the name of the output file if
you want it to be the same as the name of the input
(\filename{file.eepic} in this example).

The preferred extension for an \ePiX\ file is ``\filename{.xp}'' (for
eXtended Picture), but the source code extensions ``\filename{.c}'',
``\filename{.cc}'', ``\filename{.C}'', and ``\filename{.cpp}'' are also
permitted. If the \code{emacs} configuration file is properly set
up, \code{emacs} can recognize~``\filename{.xp}'' files, automatically
enter \ePiX\ mode, and insert a preamble template. The file
\filename{POST-INSTALL} has detailed instructions.

To process the \LaTeX\ file,
either run \LaTeX\ as usual, or do
\begin{verbatim}
laps sample
\end{verbatim}
which runs \LaTeX\ on \filename{sample.tex}, then uses \code{dvips} to
convert the \filename{dvi} to the Postscript file
\filename{sample.ps}. \code{laps} stands for ``\LaTeX\ to
Postscript''.

Figures in graphical format (as opposed to \filename{eepic}) can be
useful if you use \code{xdvi} to preview (color is available in
included files), or if you are tweaking the figures in a large
document, and do not want to recompile the entire document to change
just one figure.  \code{elaps} also handles \filename{eepic} files
(i.e., acts as \code{eepic2eps}~and \code{eepic2pdf}), even those not
produced with \ePiX.


\section{A Sample File}

Figure~\ref{fig:example} is created from the commented source file
\filename{semicirc.xp}:
% shown below. in Figure~\ref{fig:example-src}.
\begin{figure}[tb]
\begin{center}
\input{semicirc.eepic}
\caption{A semicircle.}
\label{fig:example}
\end{center}
\end{figure}
%%
%\begin{figure}[bt]
\begin{footnotesize}
\begin{verbatim}
#include "epix.h"
using namespace ePiX;                   // similar to \usepackage

// "double" = double-precision floating point
double f(double x) { return sqrt(1-x*x); } // i.e., f(x) = \sqrt{1-x^2}

double width = 0.6, height = f(width);  // rectangle dimensions
P label_here(0.5, f(0.5));              // position of graph label

int main() 
{
  unitlength("1in");
  picture(2.5, 1.25);                   // set printed size, 2.5 x 1.25in
  bounding_box(P(-1, 0), P(1, 1));      // corners; depict [-1,1] x [0,1]

  begin();   // ----- figure body starts here -----
  h_axis(4);                            // mark off 4 intervals, etc.
  v_axis(2);  
  // 2 label intervals; shift down 4pt, place below Cartesian location
  h_axis_labels(2, P(0, -4), b);    

  arc(P(0,0), 1, 0, M_PI);              // center, radius, start/end angle
  bold();                               // paths only, not fonts
  rect(P(-width,0), P(width, height));  // rectangle given by corners

  // shift label right 2pt, up 4pt, align on reference point (default)
  label(label_here, P(2,4), "$y=\\sqrt{1-x^2}$");
  end();
}
\end{verbatim}
\end{footnotesize}
%\caption{The semicircle source file.}
%\label{fig:example-src}
%\end{figure}
%%
\ePiX\ commands are of three types: definitions (of data and
functions), attribute setting, and drawing.  Like a \LaTeX\ document,
an \ePiX\ file contains a \emph{preamble}, which specifies aspects of
the figure's appearance, and a \emph{body}, which contains the actual
figure-generating commands.

Variables have been used to give logical structure to the figure: The
width of the rectangle and the the location of the label are localized
to two lines in the preamble, and the rectangle's height is computed
from the width.  The capacity to structure figures logically is one of
\ePiX's strengths. In a file as short as this, hard-wired constants
are adequate, but the importance of structuring increases rapidly with
the size of the input file, and good habits are best formed from the
start. The function definition is mostly for illustration, though does
provide slightly better organization than
\begin{footnotesize}
\begin{verbatim}
double width = 0.6, height = sqrt(1-width*width);
double lab_x = 0.5;
P label_here(lab_x, sqrt(1-lab_x*lab_x));
\end{verbatim}
\end{footnotesize}


\subsection*{Command Syntax}

Attribute-setting and drawing commands are \CXX\ \emph{functions},
blocks of instructions that can be invoked---or \emph{called}---by
name. Like mathematical functions, \CXX\ functions accept
\emph{arguments} as input and \emph{return} values.  Each argument has
a \emph{type}, such as integer (\code{int}), real number
(\code{double}, for ``double-precision floating point''), or~\code{P}
(\ePiX\ point), and may be hard-coded or specified symbolically. \CXX\ 
is a ``typed'' language, meaning that the compiler checks function
calls for the proper type and number of arguments, and issues an error
if no matching function is found.

\CXX\ allows functions to be given \emph{default arguments}; when a
function has defaults, the corresponding argument(s) may be omitted in
a call. A command is usually described by listing the types and names
of the arguments. For brevity, the type \code{double} may be left
tacit.  Default arguments are denoted by square brackets. The command
\begin{verbatim}
  line(P pt1, P pt2, [double stretch], [int n]);
\end{verbatim}
defines a ``\code{line}'' command that accepts two mandatory arguments
of type~\code{P}, an optional real-valued \code{stretch} parameter,
and an optional integer~\code{n}.  When a function is called in an
input file, arguments' types are not given: 
\begin{verbatim}
  line(P(-2,1), P(1,0), 6.4);
\end{verbatim}
The last argument is assigned its default value in this case.


\section{Basic Picture Concepts}

This section outlines the basic typographical and mathematical data
used to size and place a figure. The \emph{preamble}
of an \ePiX\ input file is everything that comes before the
\code{begin()} line; the \emph{body} comprises the portion of the file
between the \code{begin()}~and \code{end()} lines, inclusive.


\subsection*{Printed Size and Location}

\LaTeX\ treats the contents of a picture environment as a single box,
aligned by default on the lower left corner. An \ePiX\ file must tell
\LaTeX\ how large the printed figure will be, and where to align this
``picture box''. The commands
\begin{verbatim}
  picture(2.5, 1.25);
  unitlength("1in");
  offset(0.25,-0.5);
\end{verbatim}
set \LaTeX's unitlength to 1~in, create a picture environment 2.5~in
wide and 1.25~in high, then shift the picture right by 0.25~in and
down by 0.5~in. The \code{picture}~and \code{unitlength} lines are
mandatory in an \ePiX\ preamble. The \code{offset} is optional and
defaults to~$(0,0)$.\index{Offset!figure}

The argument of a \code{unitlength} command is a numerical constant
followed by one of the following valid \LaTeX\ length units: \code{bp}
(big point), \code{cm} (centimeter), \code{in} (inch), \code{mm}
(millimeter), \code{pc} (pica), \code{pt} (point, the default unit),
and~\code{sp} (scaled point).  (There are 72~big points per inch,
12~points per~pica, and 65536~scaled points per point.)  \ePiX\ does
not directly support the use of different horizontal and vertical
length units; you are expected to perform any conversions manually.

A non-zero \code{offset} causes the contents of a picture to appear in
a location where \LaTeX\ does not expect them. This may be useful when
an \code{eepic} file is included into a \LaTeX\ document, but requires
visual tweaking. A non-zero \code{offset} is risky when compiling a
figure into EPS~or PDF with \code{elaps}, since \code{dvips} may crop
the figure according to rules of its own.


\section{Logical Size, and Aspect Ratio}

\index{Aspect Ratio|(}
\index{Bounding Box|(}
An \ePiX\ figure occupies a rectangular Cartesian ``bounding box''.
The lower left and upper right corners of the bounding box are known
to \ePiX\ as \code{(x\_min,y\_min)} and \code{(x\_max,y\_max)}, while
the width and height are \code{x\_size}~and \code{y\_size}. The
bounding box is a virtual, advisory data structure; its dimensions are
not directly related to the figure's printed size, and picture
elements are not constrained to the bounding box by default.

The bounding box is specified in the preamble by giving a pair of
opposite corners; the command
\begin{verbatim}
  bounding_box(P(-1,0), P(3,2));
\end{verbatim}
sets the bounding box to be $[-1,3]\times[0,2]$. Either pair of
opposite corners may be used, though confusion is less likely if the
lower left and upper right corners are given.  Affine scaling maps the
bounding box to the picture box when the output file is written:
\begin{center}
\input{scaling.eepic}
\end{center}
\label{page:scaling}
The figure's aspect ratio is controlled by sizing the bounding box.
The aspect ratio is ``true'' if the bounding box and picture box are
geometrically similar, e.g., if both boxes are 1.5~times as wide as
they are tall.
\index{Aspect Ratio|)}


\section{Creating and Drawing Objects}

\index{Screen}
Conceptually, \ePiX\ builds a 3-dimensional virtual ``world'', then
photographs the world on a 2-dimensional ``screen''. The screen
contains the bounding box, which is affinely scaled to a \LaTeX\ 
picture environment. By default, the screen is the $(x_1,x_2)$-plane
and the world is renderd by looking straight down the $x_3$-axis.
\index{Bounding Box|)}

\subsection*{Geometric Data Structures}

The simplest object in the world is a point, represented by an ordered
triple of real numbers (double-precision floats). The
function~\code{P(x1,x2,x3)} creates the point~$(x_1,x_2,x_3)$. If only
two arguments are provided, $x_3=0$ by default. This convention allows
\ePiX\ to treat 2-~and 3-dimensional figures uniformly.  Often points
are given names, so they can be used repeatedly throughout a figure.
The (essentially equivalent) commands
\begin{verbatim}
  P pt(x_min, 2*x_min, 4);
  P pt = P(x_min, 2*x_min, 4);
\end{verbatim}
create a point named~\code{pt} and initialize it
to~$(x_\mathrm{min},2x_\mathrm{min},4)$, using the current value
of~$x_\mathrm{min}$. This value is used each time~\code{pt} occurs
subsequently.  In general, an expression involving variables may be
used to assign a value to a data structure (not just a point).
However, a variable cannot be used until a value has been assigned.
For example, the commands above will not work before the bounding box
is set, since the value of~$x_\mathrm{min}$ is unknown. Further,
changing the value of a variable does not ``update'' the values of
dependent data structures.

Points can be given in polar, cylindrical, or spherical coordinates;
all arguments are numerical.
\begin{verbatim}
  P pt=polar(r, t);  // (r*Cos(t), r*Sin(t))
  P pt=cyl(r, t, z); // (r*Cos(t), r*Sin(t), z)
  P pt=sph(r, t, phi);
\end{verbatim}
The arguments of~\code{sph} are radius (distance to the origin),
longitude (measured from the $x_1$-axis), and latitude (measured from
the equator). \index{Angular mode} By default, angles in \ePiX\ are
measured in \code{radian}s. Two other ``angular modes'' are available:
\code{degrees}~and \code{revolutions}. The angular mode is set with an
identically-named command, e.g., \code{degrees()}, and all
trigonometric operations are affected.

\ePiX\ provides algebraic operations on triples, including addition,
scalar multiplication, the dot and cross products, and a few others.
These operations are used to express relationships between triples, as
in
\begin{verbatim}
  P p1(2,1), p2(-3,1);
  P q1 = p1-p2, q2 = 2*p1-3*p2, q3=q1*q2;
\end{verbatim}
The points $q_1=p_1-p_2$, $q_2=2p_1-3p_2$, and~$q_3=q_1\times q_2$
could be given hard-coded values.  However, defining their values
symbolically imbues the figure with logical structure, making the file
easier to read, modify, and maintain. The standard basis is available:
\code{E\_1=P(1,0,0)}, etc.

Practically, ordered triples of numbers are used to represent both
\emph{locations} (points) and \emph{displacements} (vectors).
Mathematically, the concepts are distinct; for example, it makes
geometric sense to add two displacements (obtaining a displacement),
or to add a displacement to a location (to get a location), but two
locations cannot be added meaningfully. Algebraic operations act on
vectors, not points. \ePiX\ does not enforce the distinction between
points and vectors, but will in The Next Generation.


\ePiX\ implements data structures that represent line segments,
circles, planes, and spheres. These structures can be translated,
scaled, and intersected. A code snippet illustrates basic techniques:
\begin{verbatim}
  circle C1(P(0,1), P(0,-1), P(0.5,0)); // circle through 3 pts
  C1 += P(0.1,0);           // translate center
  C1 *= 2;                  // double the radius

  sphere S1(P(0,0,0), 1.5); // sphere of radius 1.5 at origin
  plane P1(P(0,0,0), E_3);  // (x,y)-plane
  circle C2 = S1*P1;        // circle of intersection
\end{verbatim}
Use of these data structures is explained in more detail in
Chapter~\ref{chapter:ref-man}.


\subsection*{Drawing}

The commands of the previous section create data structures but do not
write any output. Each type (other than~\code{P}) is drawn with
``object-oriented'' syntax.  For example, if~\code{C1} is a circle,
then the command \code{C1.draw()} draws the circle. The effect of
drawing a \code{plane}~or \code{sphere} is described in
Chapter~\ref{chapter:ref-man}. \ePiX\ also provides high-level
commands that draw polygons, curves, and compound objects, such as
arrows\index{Arrow} and coordinate axes.  Drawing commands must come
in the body of the file, after \code{begin()}.

\begin{verbatim}
  line(P p1, P p2);
  triangle(P p1, P p2, P p3);   // vertices specified
  rect(P p1, P p2);             // coord rect w/opposite corners
  quad(P p1, P p2, P p3, P p4); // arbitrary quadrilateral

  spline(P p1, P p2, P p3);     // quad/cubic splines given 
  spline(P p1, P p2, P p3, P p4); // by control points

  arc(P center, radius, t_min, t_max);
  ellipse(P ctr, P v1, P v2, [t_min], [t_max], [int n]);
\end{verbatim}
%  polyline(int n, P p1, ..., P pn);
%  polygon (int n, P p1, ..., P pn);
The arguments of \code{rect} must lie in a plane parallel to a
coordinate plane; the sides of the rectangle are parallel to
coordinate axes.

An \code{arc} has the given center and radius, and lies in a plane
parallel to the $(x_1,x_2)$-plane. Angles are measured from the
\code{E1}~direction in the current angle units.

An \code{ellipse} draws an elliptical arc with the specified center
and ``axes''; precisely, the curve drawn is parametrized by
\[
t\mapsto \mathrm{ctr} + (\cos t)v_1 + (\sin t)v_2,\qquad
t_\mathrm{min}\leq t\leq t_\mathrm{max}.
\]
(Note that \code{ctr} is a location, while \code{v1}~and \code{v2} are
displacements.)  As with \code{arc}s, angles are measured in current
angle units. If the parameter bounds are omitted, the entire ellipse
is drawn. The final (optional) argument says how many points to use
when drawing the arc. This can be omitted safely in most situations.


\section{The Camera}

\index{Camera|(}
\index{Screen|(}
Art students sometimes practice perspective by tracing on a window
with grease pencil, a mathematical transformation called \emph{point
projection}. \ePiX's default mapping from the world to the screen is
similar. Imagine standing before a scene at a \emph{viewpoint}, and
possessing X-ray vision, so that objects are transparent. Somewhere in
front of you is a plane, the screen. The \emph{target} is the point on
the screen obtained by dropping a perpendicular from the viewpoint.

Three mutually perpendicular unit vectors sit at the target: sea, sky,
and eye. The sea vector points to the right, sky points upward, and
eye points straight from the target to the viewpoint. The bounding
box\index{Bounding Box} of the figure is defined with respect to the
screen's (Cartesian) sea-sky coordinate system.

Given a point~$p$ in front of the viewpoint, we want to determine the
screen location to which~$p$ projects. Join~$p$ to the viewpoint by a
line; this line intersects the screen plane exactly once, and the
point of intersection is where we draw~$p$ in the screen. 
\begin{figure}[hbt]
  \begin{center}
    \input{camera.eepic}
  \end{center}
  \caption{Point projection.}
  \label{fig:camera}
\end{figure}
\index{Screen|)}

\index{Camera!manipulation}
At the start of a figure, the camera is initialized to lie on the
$x_3$-axis at very large distance from the origin. The resulting view,
essentially projection along the axis, is suitable for 2-dimensional
figures. The camera is manipulated with object-oriented syntax:
\begin{verbatim}
  camera.at(P posn);               // set viewpoint to posn
  camera.look_at(P targ);          // set target
  camera.range(double dist);       // fix target, move viewpoint
  camera.focus(double dist);       // fix viewpoint, move target
  camera.rotate_sea(double angle); // rotation about an axis
\end{verbatim}
These commands must come in the figure body.
\index{Camera|)}


\section{Drawing Attributes}

\ePiX\ creates line drawings, not pixmaps.  In an \code{eepic} file,
objects are either pointlike (\LaTeX\ glyphs, text boxes) or pathlike
(everything else). The appearance of a pathlike object depends on the
line style, width, filling,\index{Filling} and color. Each attribute
is a declaration, remaining in effect until superceded or the file
ends.  When a file starts, paths are drawn solid, black, unfilled, and
with \code{thinlines} (a width of $0.4$~pt) in effect.
Attribute-setting commands should come in the figure body.


\subsection*{Path Width and Style}

\index{Path!width}
\index{Path!style|(}
The standard path widths are \code{plain}~(\code{thinlines}) and
\code{bold}~(\code{thicklines}). Other widths are available via a
\code{pen()} command. The argument is a number, interpreted as a
length in~\code{pt}, or a number followed by a two-letter \LaTeX\ 
length as in the \code{unitlength()} command.  Journals discourage
line widths smaller than about 0.5~pt, and a multitude of
author-specified line widths tends to look cluttered and \emph{ad
  hoc}. If possible, use the standard widths.
\begin{verbatim}
  plain();        // thinlines, about the same as pen(0.4);
  bold();         // thicklines,                  pen(0.8);
  pen("0.02in");  // set path width to 0.02 in
\end{verbatim}

The path \emph{style} is one of \code{solid}, \code{dashed}, or
\code{dotted}. The style is set with an identically-named command,
e.g., \code{dashed()}.  Every pathlike object in \ePiX\ is drawn using
a list of points. When the path style is \code{solid}, the points are
joined connect-the-dots fashion. A \code{dotted} path is drawn by
placing a small dot at each point of the path. When the style is
\code{dashed}, line segments are drawn partway from each vertex to its
neighbors.
\begin{center}
  \input{pathstyle.eepic}
\end{center}
A few parameters can be adjusted manually: the page distance between
consecutive points on a path (for polygons---triangles, quads, etc.),
the ``dash density'' (the percentage of a dashed path filled by
dashes), and the dot size:
\begin{verbatim}
  dash_fill(0.7); // dashes fill 70% of point gap
  dash_length(6); // path points separated by 6pt
  dot_sep(8);     // path points separated by 8pt
  dot_size(2);    // dots 2pt in diameter
\end{verbatim}
\index{Path!style|)}


\subsection*{Color and Filling}

\index{Color|(}\index{PSTricks}
An \ePiX\ output file achieves color through the \code{pstcol} style,
an amalgamation of \code{color}~and \code{PSTricks} styles. There are
three available color models: \code{rgb}, \code{cmyk} (cyan, magenta,
yellow, and black), and named (primary colors only). An \code{rgb}
color is specified by three numbers between~0 (no color) and~1 (full
saturation), each of which represents the density of a primary color
(red, green, and blue respectively). A \code{cmyk} color is similarly
set by giving four densities. Finally, a primary color (r, g, or~b; c,
m, or~y; black and white) may be specified by name, with an optional
density.
\begin{verbatim}
  rgb(1, 0, 1);  // magenta
  cmyk(0,1,0,0); // same thing
  magenta();     // method III
  rgb(1, 0.7, 0.7); // light red
  rgb(0.4, 0, 0);   // dark red
  red(0.4);         // same thing
\end{verbatim}
Like all parameters, color densities can depend on variables. Values
that lie outside the interval~$[0,1]$ are ``clipped''; for example,
\code{rgb(1.4,-0.05,2)} is also magenta.
\index{Color|)}

\index{Filling}\index{Path}
The \code{fill()} command causes closed paths to be gray-shaded, using
Post\-Script \code{special}s. The depth of gray ranges from~0 (white)
to~1 (black), and defaults to~$0.3$. The command \code{gray(0.4)} sets
the depth to~$0.4$.  Filling is deactivated by the command
\code{fill(false)}. The order of filled objects in the source file is
significant, because \ePiX\ writes its output in the same order, and
PostScript has no transparency. Layering, hidden object removal, and
color shading are discussed in Chapters \ref{chapter:ref-man}~and
\ref{chapter:adv}.


\section{Typography}

\index{Labels|(emph}
In an \ePiX\ file, points are denoted with \LaTeX\ glyphs
(``markers'') or text boxes (``labels''). A marker occupies a box of
zero size, and is placed at a specified Cartesian location. A label
has typographical size, so its placement is more involved.  When the
size or aspect ratio\index{Aspect Ratio} of a figure is adjusted, the
font size stays the same. In order to keep a label aligned properly
over a range of sizes, a scale-invariant alignment point is attached
to each label, and Cartesian coordinates are used to position the
alignment point. The alignment point is controlled both by a true-pt
offset and an optional \LaTeX-style alignment option.


\subsection*{Markers}

\index{Marker types|(emph} \ePiX\ markers are called
by name: \input{dots.eepic}
\begin{verbatim}
  spot(P pt);  dot(P pt);   ddot(P pt);
               box(P pt);   bbox(P pt);  
  ring(P pt);  circ(P pt);
\end{verbatim}
\code{spot} and \code{(d)dot} are solid dots that cannot be colored.
\code{box}~and \code{bbox} are solid, colorable squares.  A
\code{ring} is colorable and transparent, while a \code{circ} is
uncolorable and opaque. Markers in the same column are the same size,
and each column is~1.5 times the diameter of the next. The diameter of
a \code{dot} (hence that of all the above \code{marker}s) is set with
\code{dot\_size(diam)}. The argument, 3~by default, is a number
of~\code{pt}.  The glyphs listed in Table~\ref{table:marker} are
available via the command
\begin{verbatim}
  marker(P pt, <MARKER TYPE>);
\end{verbatim}
\begin{table}[tb]
  \begin{center}
    \input{marker.eepic}
  \end{center}
  \caption{\ePiX's \code{marker} types.}
  \label{table:marker}
\end{table}
%% These marker types are also available when plotting data from a file.
\index{Marker types|)}


\subsection*{Labels}

A \emph{label} is a typographical box. Since a \LaTeX\ box occupies a
rectangle on the printed page, a single location is not enough
information to position a label within a figure; an alignment point is
needed in addition.  By default, the alignment point of a text box is
its reference point, the intersection of the left edge and the
baseline, which is used by \LaTeX\ to position the box on the
page:\input{basepoint.eepic}
An alignment point may be shifted manually:
\begin{verbatim}
  label(P(3,2), P(2,-1), "$\\rho=\\sin \\theta$");
\end{verbatim}
typesets the equation $\rho=\sin\theta$ and places the resulting text
box at Cartesian location~$(3,2)$, but shifted right by~\code{2pt} and
down by~\code{1pt}. Note that \CXX~treats ``\verb+\\+'' as a
``backslash control sequence'', so a double backslash is needed in the
source to get a single backslash in the output.  The general
\code{label} commands are:
\begin{verbatim}
  label(P posn, P offset, <label text>);
  label(P posn, <label text>);
  label(P posn, P offset, <label text>, align);
\end{verbatim}
\index{Offset!label|(}
The first command prints a label with its \LaTeX\ reference point at
the Cartesian location~\code{posn}, offset in true points (on the
page) by the specified amount. Offsets of 2, 4, 6, 12, or~18 points
work well with a 12-point font.\index{Labels!offset}

The second command prints the label text in a \LaTeX\ box centered
at~\code{posn}. While this is perhaps the most obvious way of placing
a label, it may not be the correct method, since labels often mark a
geometric object that should not be covered by the label.

\index{Labels!alignment|(}
In the third command, the \code{align} option may be one---or an
appropriate pair---of \code{t}, \code{b}, \code{r}, or~\code{l} (top,
bottom, right, left), or~\code{c} (center). As with \code{offset}s,
these alignment options specify the position of the label
\emph{relative to the Cartesian location}~$p$, namely they work
\emph{opposite} to the way they work in \LaTeX. For example, the
alignment option~\code{br} puts the label below and to the right
of~$p$.
\begin{center}
  \input{alignment_lr.eepic}
  \input{alignment.eepic}
\end{center}
\index{Labels!alignment|)}

Each label command has a corresponding ``mask'' version
(\code{masklabel}) that draws an opaque white rectangle under the
label text. Masking is useful when the text of a label sits in a
cluttered part of the figure. An \eepic\ file containing
\code{masklabel}s requires the \code{pstcol} package.

Labels can be rotated; the angle is set in current angle units with
the command \code{label\_angle(theta)}. A rotation angle
of~90 degrees prints labels suitable for a vertical axis.  An
\eepic\ file containing rotated labels requires the \code{rotating}
package.

When constructing and placing a label, keep in
mind that
\begin{itemize}
\item Alignment offsets are specified in \code{pt} (i.e., page
  coordinates), not in Cartesian units, because the alignment point
  should not depend on the logical or printed size of the figure.

\item The label text is enclosed in double quotes (the single
  character~\code{"}), and contains the \LaTeX\ code to generate the
  label. Backslashes are doubled.
\end{itemize}
\index{Offset!label|)}
\index{Labels|)}


\section{\CXX\ Basics}

An \ePiX\ source file is a \CXX~program. If you've successfully
modified and compiled any of the sample files, you know enough~\CXX\ 
to use \ePiX. In the author's experience, \code{C}~grammar suffices
for most applications.  An excellent introduction to definitions of
functions and variables, control statements, and overall program
structure is Kernighan~and Ritchie's \emph{The \code{C} Programming
  Language}, second edition~\cite{KnR}.


Jay Belanger's \code{emacs} mode for \code{ePiX} inserts a file
template when an empty buffer is opened with the extension~\code{xp}.
This section explains the purposes served by the template. A few
additional remarks may help you avoid basic syntax pitfalls.


\subsection*{Comments}

\CXX\ has two types of comments. \code{C}-style comments, which may
span several lines, are delimited by the strings \code{/*}~and
\code{*/}. One-line comments, analogous to the \LaTeX~\code{\%}, are
begun with \code{//}. A one-line comment may appear within a multiline
comment, but a \code{C}-style comment may not; the compiler will
mistake the first~\code{*/} it encounters as the end of the current
multiline comment.


\subsection*{Pre-Processing}

The compiler ignores nearly all whitespace (spaces, tabs, and
newlines), which should be used liberally to make files easy to read.
Other punctuation (periods, commas, (semi)colons, parentheses, braces,
and quotes) dictates file parsing, and must adhere stringently to
grammar.

A \CXX\ file consists of ``statements''. A statement ends with a
semicolon, and conventionally a file contains one statement per line
(when possible). For historical reasons, an input file is
``pre-processed'' before the compiler sees it. The most important use
of pre-processing in \ePiX\ is file inclusion.  An \ePiX\ file always
begins with the lines
\begin{verbatim}
#include "epix.h" // N.B. pre-processor directive, no semicolon
using namespace ePiX;
\end{verbatim}
The first line is analogous to a \LaTeX\ \code{usepackage} command: It
causes the pre-processor to replace the line with the contants of a
file, thereby importing the names of commands provided by~\ePiX.  To
avoid name conflicts, \ePiX's commands are enclosed in a
``namespace''. For example, the \code{label} command is actually known
to the compiler as \code{ePiX::label}. The second line above tells the
compiler to apply the prefix tacitly.

\subsection*{Variables and Functions}

Definitions of variables and functions play the same role in a figure
that macro definitions do in a \LaTeX\ document: gathering and
organizing information on which the figure depends. A variable is
defined by supplying its type, name, and initial value. By far the
most common data types in~\ePiX\ are \code{double}, \code{P}, and
\code{int}. The name of a variable may consist (only) of letters
(including the underscore character) and digits, and must begin with a
letter:
\begin{verbatim}
my_var, my_var2,  _MY_var, __, aLongVariableName;  // valid
my-var,    2var, v@riable, $x, ${MY_VARIABLE}; // not valid
\end{verbatim}
Variable names are case-sensitive. There are numerous conventions
regarding the significance of capitalization. Generally, make names
descriptive but not unwieldy, and avoid names that begin with an
underscore (unless you know what you're doing).

A function accepts ``arguments'' and ``returns a value''. To define a
function in~\CXX, you must specify the return type, the name of the
function, the types of the arguments, and the algorithm by which the
value is computed from the inputs. The code block
\begin{footnotesize}
\begin{verbatim}
double f(double x) 
{ 
  return sqrt(1-x*x); 
}
\end{verbatim}
\end{footnotesize}
specifies the \code{double}-valued function~$f$ of one \code{double}
variable defined by the formula $f(x)=\sqrt{1-x^2}$. Several sample
and source files (especially \filename{functions.cc}) give more
interesting examples. A function definition should be formatted as
above for readability.


\subsection*{Program Execution}

All the ``action'' in a \CXX\ program occurs inside the special
function~\code{main}. Running a compiled \CXX~program is viewed by the
operating system as calling the program's \code{main} function. The
return value (an \code{int}) is the program's exit status.  In an
\ePiX\ file, the \code{main} action usually consists of setting the
logical and printed size of the figure, then building and drawing the
figure, changing attributes when desired.  \ePiX\ output proper starts
with \code{begin()} and terminates with \code{end()}. The intervening
statements constitute the \emph{body} of the figure.

In \CXX, a function may not be defined inside another function. Thus,
variables may be defined inside \code{main}, but functions cannot be.


\subsection*{Raw output} 

More-or-less verbatim text can be printed to the output file.  A
single backslash is produced by a double backslash in the input file.
Certain letters have special meanings when backslash-escaped,
including ``\verb+\n+'' (newline) and ``\verb+\t+'' (TAB).
Unlike~\LaTeX, \CXX~does not require a space to separate an escape
sequence from following text; the string ``\verb+\textwidth+'' is read
``\code{TABextwidth}'' by the compiler.

As an application, a complete \LaTeX\ \code{figure} environment (with
caption and label) can be produced by an \ePiX\ file. Newlines must be
added explicitly, and all printing statements must occur inside the
call to \code{main()}.
\begin{verbatim}
#include "epix.h"
using namespace ePiX;
using std::cout;      // C++'s output function

int main() 
{
  cout << "\\begin{figure}[hbt]\n";
  unitlength(...);    // picture, bounding_box, etc.
  begin();
  < ... ePiX commands ... >
  end();
  cout << "\\caption{A \\LaTeX\\ figure.}\n" // N.B. line con-
       << "\\label{fig:example}\n"           // tinues, no ";"
       << "\\end{figure}\n%%%%\n";           // LaTeX formatting
} // End of main()
\end{verbatim}


\subsection*{Conditionals and Loops}

An algorithm's behavior usually depends on some internal state. A
\emph{conditional statement} causes blocks of code to be executed
depending on some criterion. A \emph{loop} repeatedly executes a code
block, usually changing the values of variables in a predictable way,
so that the loop exits after finitely many traversals.
Figure~\ref{fig:gcd} illustrates both conditionals and loops with
Euclid's algorithm for the greatest common divisor. Three pieces of
notation require explanation: \code{j\%i} means~``$j\pmod i$'',
\code{||} is logical~``or'', and \code{==} is ``test for equality''.
(A single~``\code{=}'' is the assignment operator.)
\begin{figure}[hbt]
\begin{footnotesize}
\begin{verbatim}
int gcd (unsigned int i, unsigned int j)
{
    int temp=i;
    if (i==0 || j==0)
        return i+j;    // define gcd(k,0) = k

    else {
        if (j < i)     // swap them
            {
                temp = j;
                j=i;
                i=temp;
            }
        // the work is done here...
        while (0 != (temp = j%i)) // assign temp, test for zero
            {
                j=i;
                i=temp;
            }
        return i;
    }
}
\end{verbatim}
\end{footnotesize}
\caption{Euclid's division algorithm in~\CXX.}
\label{fig:gcd}
\end{figure}

%% \noindent For details on conditionals, or for more advanced techniques
%% of \CXX\ programming, please consult a textbook or online tutorial.


\section{High-Level Picture Elements}

\ePiX\ implements a miscellany of high-level drawing capabilities:
arrows, coordinate axes and axis labels, polar plots, data plotting
from files, calculus operations, vector fields, solutions of
differential equations, and non-Euclidean geometry.


\subsection*{Arrows} 

\index{Arrow|(emph}
An arrow is specified by its tail and tip. An optional third argument
scales the arrowhead.
\begin{verbatim}
  arrow(P tail, P tip, [double scale]);
  dart (P p1, P p2);  // same as arrow(p1, p2, 0.5);
  aarrow(P p1, P p2); // double-headed arrow <--->
\end{verbatim}

Pictorially, an \code{arrow} consists of a line segment (the shaft)
surmounted by a triangle (the arrowhead). In profile, an arrowhead's
width is \code{3pt}, and its height is 5.5~times the width.  The
actual printed height depends on the \code{arrow}'s orientation with
respect to the camera.  

By default, an arrowhead is a hollow triangle, which can be colored.
The \code{fill()}\index{Filling} command produces solid, uncolorable
arrowheads.  \code{PSTricks}\index{PSTricks} commands can be used for
solid colored \code{arrow}s.  \index{Arrow|)}


\subsection*{Coordinate Axes and Labels} 

\index{Axes|(}
A coordinate axis consists of a line between two points together with
a specified number of regularly-spaced tick marks:
\begin{verbatim}
  h_axis(p1, p2, n);    // n subintervals (n+1 ticks)
  v_axis(p1, p2, n);
\end{verbatim}
The style of tick mark is appropriate for an axis of the given type.
If the endpoints are omitted, they default to $p_1=(x_\mathrm{min},0)$
and $p_2=(x_\mathrm{max},0)$ for a horizontal axis, or to
$p_1=(0,y_\mathrm{min})$ and $p_2=(0,y_\mathrm{max})$ for a vertical
axis. If the bounding box\index{Bounding Box} has integer width and/or
height, then omitting the number of points draws tick marks one unit
apart.

\index{Labels!axis|(}
Labels for a horizontal axis are generated with:
%%  h_axis_labels(p1, p2, n, P(u,v));
\begin{verbatim}
  h_axis_labels(P p1, P p2, int n, P offset, [alignment]);
\end{verbatim}
This puts $(n+1)$~evenly-spaced labels on the segment joining
\code{p1}~and \code{p2}. The labels are automatically generated to
match their horizontal location. As for ordinary labels, the
\code{offset} is in~\code{pt}, and the optional \LaTeX-style alignment
option places the labels using their corners rather than their
reference points. Labels for a vertical axis are generated in
the obvious way.\index{Offset!label}

As for coordinate axes, the initial and final points may be omitted in
an \code{axis\_label} command, with the same defaults.  However, the
\code{offset} and number of labels must always be specified.
\index{Labels!axis|)}


\subsection*{Coordinate Grids} 

Cartesian grids fill a coordinate rectangle, and have a specified
number of lines in each direction. A polar grid has specified radius,
rings, and sectors.
\begin{verbatim}
  grid(n1, n2);         // fills the bounding box
  grid(p1, p2, n1, n2); // fills the box with corners p1, p2
  grid(p1, p2, mesh(n1, n2), mesh(m1,m2));
  polar_grid(r, n1, n2);
\end{verbatim}
Each command draws an \code{n1}~by \code{n2} grid. The third uses
an $m_1\times m_2$ mesh, which is useful only if the camera lens does
not map lines in object space to lines on the screen.

\index{Graph paper|(}
Graph paper may be created by superimposing grids:
\begin{center}
  \begin{minipage}[b]{2.5in}
    \begin{footnotesize}
\begin{verbatim}
  pen(0.25);
  grid(10*x_size, 10*y_size);
  pen(0.5);
  grid(2*x_size, 2*y_size);
  pen(1);
  grid(x_size, y_size);
\end{verbatim}
    \end{footnotesize}
  \end{minipage}
  \qquad
  \input{graphpaper.eepic}
\end{center}
  
S. D. Pedersen's \code{contrib/} package provides enhanced Cartesian
graphs. See \code{contrib/doc} in the source for documentation.
\index{Graph paper|)}
\index{Axes|)}


\section{Basic Plotting}

\index{Plotting|(}
Because \filename{eepic.sty} can draw lines of arbitrary length and
slope, curves can be approximated by connect-the-dots paths. \ePiX\ 
renders curves, polygons, and function graphs in this way.

For the moment, ``function'' means ``function of one variable''
(precisely, a \code{double}-valued function of a \code{double}
variable). A function graph depends on the domain and the number of
points to use. Each of the commands
\begin{verbatim}
  plot(f, t_min, t_max, n);
  polarplot(f, t_min, t_max, n);
  shadeplot(f, t_min, t_max, n);
\end{verbatim}
graphs the function~\code{f} on the interval \code{[t\_min, t\_max]}
by dividing the interval into \code{n}~subintervals of equal length.
The first gives a Cartesian plot, the second a polar
plot\index{Plotting!polar} with bounds in current angular units, the
third shades the region between the graph and the horizontal axis. If
two functions are given to \code{shadeplot}, the region between their
graphs is shaded.


\subsection*{Data plotting} 

\index{Plotting!data|(}
Files of numerical data can be manipulated, analyzed, and plotted.
The format for a data file is one or more floating-point numbers per
line, with the same number of entries per line. A line in a data file
that starts with a~\code{\%} is a comment. 

Roughly, the basic plot command reads numbers from two (or three)
columns of a specified file, treats them as coordinates, and plots the
resulting points.
\begin{verbatim}
  plot("filename", STYLE, columns, [i_1], [i_2], [i_3], [F]);
\end{verbatim}
The first argument is the name of the data file. The \code{STYLE} may
be \code{PATH}, which joins the points in the order they appear, or
any of the marker types\index{Marker types} in
Table~\ref{table:marker}.  Next comes the number of columns (entries
per line); if there are fewer columns than expected in the file,
nothing is plotted, while if there are more columns than expected, a
warning is issued but the data is plotted.  The remaining arguments
are optional. The integers~$i_k$ specify columns from which to extract
data. The column entries are fed to the \code{P}-valued
function~\code{F} to obtain points. If~\code{F} is omitted, it
defaults to the Cartesian point constructor. If~$i_3$ is omitted, the
$i_1$~and $i_2$~columns are used to create points; by default,
$i_1=1$~and $i_2=2$. For example, if \filename{mydata.dat} contains
6~columns of numbers, then the respective commands
\begin{verbatim}
  plot("mydata.dat", DOWN, 6);
  plot("mydata.dat", BOX, 6, 2, 4, 5, sph);
\end{verbatim}
plot the first two columns of \filename{mydata.dat}, putting
a~``$\bigtriangledown$'' at each point; and extract the second,
fourth, and fifth columns of the file, treat them as spherical
coordinates, and put a~``\input{box.eepic}'' at each point.

For more elaborate (e.g., user-defined) analysis, data may be read
into a \code{FILEDATA} structure. A \code{FILEDATA} is a
\CXX~vector of columns, each column having as many entries as there
are lines of data in the file. The snippet below reads data from a
file, then plots the result.
\begin{verbatim}
  FILEDATA my_cols(6);          // vector of 6 columns
  read("mydata.dat", my_cols);  // read data from file
  ...                           // code to mangle data
  plot(my_cols, BOX, 2, 4, 5, sph);
\end{verbatim}
The number of columns needn't be specified in the \code{plot} command,
since it was provided when the data was read.  The post-\code{TYPE}
options are identical to the earlier plot command.  The $j$th~entry of
the $i$th~column is called \code{my\_cols.at(i).at(j)}.

\index{Linear regression}
\ePiX\ implements simple numerical functions of a \code{FILEDATA}
structure:
\begin{verbatim}
  avg(my_cols, i_1);     // arithmetic mean of the i_1 column
  var(my_cols, i_1);     // variance of the i_1 column
  covar(my_cols, i_1, i_2); 
  regression(my_cols, i_1, i_2);  // draw regression line
\end{verbatim}
The covariance of two columns is obtained by subtracting the
respective averages entry by entry, then taking the dot product. The
regression line is the least-squares best fit for predicting
column~$i_2$ from column~$i_1$.
\index{Plotting!data|)}
\index{Plotting|)}


\chapter{Reference Manual}
\label{chapter:ref-man}

This chapter covers the design and use of \ePiX, assuming you've
thoroughly digested the material in Chapter~\ref{chapter:started}.
Remaining features are documented, and the implementation described.
If a feature isn't explained here, please consult the source code or
contact the author.


\section{More About \CXX}

A textbook or similarly detailed reference is essential for serious
study of~\code{C} or~\CXX. Go with the masters: \emph{The \code{C}
  Programming Language}, second edition, by Brian Kernighan~and Dennis
Ritchie~\cite{KnR}, is an excellent, manageable resource for the
basics of procedural programming, while \emph{The \CXX~Programming
  Language}, by Bjarne Stroustrup~\cite{S}, is a definitive,
encyclopedaic reference.

\CXX\ is a powerful, complex language whose syntax is similar to that
of~\C, or to the scripting languages of Maple~and Mathematica. An
\ePiX\ input file is source code for a \CXX~program that writes an
\eepic\ file as output.  \ePiX\ may be viewed as an extension to~\CXX;
in the same way that \LaTeX\ furnishes a high-level interface to~\TeX,
\ePiX\ provides a high-level bridge between the computational power
of~\CXX\ and the \LaTeX\ \code{picture} environment. 

Like all high-level programming languages, \CXX\ provides variables,
functions, and control structures. Variables hold pieces of data such
as numerical values and geometric locations, while functions operate
on data. A control structure, such as a loop or conditional statement,
affects the program's course according to the program's current state.
A source file is composed primarily of ``statements'', which perform
actions ranging from defining variables and functions to setting
figure attributes, performing calculations, and writing objects to the
output file.

\subsection*{Names and Types}

Names of variables and functions may consist (only) of letters,
digits, and the underscore character. The first character of a name
must not be a digit, and the language standard reserves names starting
with underscore for library authors. Names are case-sensitive, but
it's usually a bad idea to use a single name capitalized and
uncapitalized in a single file. Numerous capitalization conventions
are used informally; this document uses uncapitalized words separated
by underscores for variables and functions, and occasionally uses all
capitals for constants. As with names of \LaTeX\ macros, primary
considerations are clarity (of meaning), readability, and consistency.

Every variable in \CXX\ has a ``type'',
such as integer (\code{int}), double-precision floating point
(\code{double}), or Boolean (\code{bool}, true~or false).  \ePiX\ 
provides additional types, the most common of which is~\code{P}, for
point. The construct \code{P(x,y,z)} creates~$(x,y,z)$, while
\code{P(x,y)} gives~$(x,y,0)$, which is effectively the pair~$(x,y)$.
A variable is defined by giving its type, its name, and an
initializing expression.

In~\code{C} and~\CXX, a \emph{pointer} is a variable that holds the
memory address of another variable. Though more subtle than ordinary
variables, pointers are useful in expressing certain algorithms, such
as sorting. In Japan, buildings' addresses are assigned
chronologically, rather than according to street location. A building
is analogous to a variable, while the address is a pointer. If the
Japanese parliament passed a law mandating that buildings be addressed
consecutively along the street, there would be two ways to proceed:
Dig up and relocate each physical building (move variables), or
re-number the buildings in place (sort pointers). For similar reasons
of efficiency, \CXX's sorting algorithms work with pointers.

\CXX\ also provides \emph{reference} variables, which allow a variable
to be given an additional name. Their use arises because of the way
\CXX\ functions treat their arguments.


\subsection*{Functions}

\index{Functions|(} In a programming language, the term ``function''
refers to a block of code that is executable by name.  A \CXX\ 
function takes a list of ``arguments'', and has a ``return value''.
This information, together with the function's name, must be provided
when a function is defined.  A function may not be defined inside
another function. However, a function may call other functions
(including itself) as part of its execution:
\begin{verbatim}
  int factorial(unsigned int n)
  {
    if (n == 0) return 1;
    else return n*factorial(n-1);
  }
\end{verbatim}

The special type~\code{void} represents a ``null type''. A function
that performs an action but does not return a value has return type
\code{void}. A function that takes no arguments may be viewed as
taking a single \code{void} argument.

Every \CXX\ program has a special function~\code{main()}, which is
called by the operating system when the program is run. The arguments
of~\code{main()} are command-line arguments, and the return type is an
integer that signals success or failure. User-specified functions must
be defined before the call to \code{main()} or in a
separately-compiled file.

Functions in~\CXX\ may be as simple as an algebraic formula or as
complex as an arbitrary algorithm. Greatest common divisors, finite
sums, numerical derivatives and integrals, solutions of differential
equations, recursively generated fractal curves, and curves of best
fit are a few applications in \ePiX. Several sample files contain
user-level algorithms, which do not require knowledge of \ePiX's
internal data structures. The source file \filename{functions.cc}
contains simple functions defined by algorithms, and
\filename{functions.h} illustrates the use of \CXX\ templates. Other
source files, such as \filename{plots.cc}, may be consulted for
Simpson's rule, Euler's method, and the like.

An error, such as division by zero or an attempt to intersect parallel
lines, may occur when a function is executed. In this situation, a
\CXX\ function can ``throw an exception'', or return an error type
that the caller ``catches'' and handles. If an uncaught exception is
thrown, the program terminates. \ePiX's intersection operators throw
exceptions when certain conditions are not met.


\subsection*{Mathematical Functions}

\CXX~knows several familiar mathematical functions by name:
\begin{verbatim}
  sqrt   exp   log   log10   ceil   floor   fabs
\end{verbatim}
(\code{fabs} is the absolute value for a floating-point argument.)
\ePiX\ provides trig functions and their inverses that are sensitive
to angular mode:\index{Angular mode}
\begin{center}
\begin{tabular}{p{1.25in}p{1.25in}p{1.25in}}
  \code{Cos}  & \code{Sin}  & \code{Tan} \\
  \code{Sec}  & \code{Csc}  & \code{Cot} \\
  \code{Acos} & \code{Asin} & \code{Atan}
\end{tabular}
\end{center}
The inverse trig functions are principle branches.

The function \code{pow(x,y)} returns~$x^y$ when $x>0$, and
\code{atan2(y,x)} (N.B. argument order) returns
$\mathrm{Arg}(x+iy)\in(-\pi,\pi]$, the principle branch of arg.
\CXX~knows many constants to 20~decimal places, such as \code{M\_PI},
\code{M\_PI\_2}, and \code{M\_E} for~$\pi$, $\pi/2$, and~$e$
respectively.  \ePiX\ defines a few additional functions:
\begin{verbatim}
  recip   sgn   zero   sinx   cb   id   proj1   proj2
\end{verbatim}
\begin{center}
  \input{sgn.eepic}\qquad
  \input{sinx.eepic}\qquad
  \input{cb.eepic}
\end{center}
\code{recip} is the reciprocal, defined to be~$0$ at~$0$; \code{sgn}
is the signum function; \code{zero} is the constant function;
\code{sinx} is the function $x\mapsto \sin(x)/x$ with the
discontinuity removed; \code{cb} (for ``Charlie Brown'') is the
period-2 extension of the absolute value function on~$[-1,1]$;
\code{id} is the identity mapping, defined for an arbitrary data type;
the \code{proj} functions return their first and second variable,
regardless of type.

The GNU \CXX\ library defines other functions, including inverse
hyperbolic functions (\code{acosh}, etc.), \code{log}~and \code{exp}
with base~2, 10, or arbitrary~$b$ (\code{log2}, etc.), the error and
gamma functions (\code{erf}~and \code{tgamma} [sic], respectively),
and Bessel functions of first and second kind: \code{j0}, \code{j1},
\code{y0}, etc. Use, e.g., \code{jn(5,\ )} to get higher indices.
The GNU \C~library reference manual~\cite{GNUC} describes these and other
functions in detail.

Functions may be used in subsequent definitions, and functions of two
(or more) variables are defined in direct analogy to functions of one
variable:
\begin{verbatim}
double f(double t) { return t*t*log(t*t); } // t^2 \ln(t^2)
double g(double s, double t) { return exp(2*s)*Sin(t); }
\end{verbatim}


\subsection*{Basics of Classes}

Unlike~\code{C}, \CXX\ supports ``object-oriented programming''. In a
nutshell, a \emph{class} is an abstraction in computer code of some
concept, such as a point, a sphere, a mapping that can be plotted, or
a camera. Implementationally, a class consists of \emph{members}
(named data elements) and \emph{member functions} (functions that
belong to the class and have free access to members).  The idea is to
encapsulate objects and the operations that make sense for them in a
single logical entity.  In simple programming, classes may be treated
like built-in types.

\CXX\ classes enforce access permissions on their members, protecting
data from being manipulated except in controlled ways. A class's
member functions (and ``public'' members) constitute its \emph{user
  interface}. The concept of class separates logical aspects of a data
type from a particular implementation.

Each ``instantiation'' of a class has its own member functions.
A member function therefore knows ``which object'' it was called on,
and the call syntax differs from standard function calls:
\begin{verbatim}
  circle C1(P(1,0), 1.5); // circle of given center and radius
  C1.draw();              // member function circle::draw();
\end{verbatim}
Naturally, this call draws the circle~\code{C1}.

A few short paragraphs cannot do more than scratch the surface of
classes and object-oriented programming. Please consult a book, such
as Stroustrup~\cite{S}, for details.


\subsection*{References and Function Arguments}

\code{C}~and \CXX\ are ``call by value'' languages. Actual variables
are not passed to a function; instead a copy of the value is made, and
the function operates only on the copy. Though this feature causes
occasional inconvenience, it prevents a variable from having its value
changed by a function.

A function may accept arbitrary data structures as arguments. If a
data structure is complicated and a function is called frequently
(thousands or millions of times), implicit copying becomes a source of
inefficiency. References bypass this problem: If a variable is passed
by reference, only a pointer need be copied. 

A variable that is passed by reference may be altered by the calling
function. This technique of updating variables is often touted as a
feature in \CXX\ texts; however, such trickery circumvents the data
encapsulation of calling by value. Passing a variable by reference so
that a function can modify the value of a variable is considered less
than ideal programming practice.


\subsection*{Overloading}

\CXX\ provides ``overloading'': Multiple functions can be given the
same name, so long as the number and/or type of their arguments
differ. (It is \emph{not} enough for the return types alone to differ.
The compiler must be able to select a function from its calling
syntax.)  To the user, the appearance is that a single function
intelligently handles multiple argument lists. Naturally, overloaded
names should refer to functions that are conceptually related.
Overloading tends to be most useful in library code; \ePiX\ provides
numerous \code{plot} functions, for instance.
\index{Functions|)}


\subsection*{Scope}

A \CXX\ statement ends with a semicolon. A collection of statements
enclosed by curly braces is a ``code block'', and may be viewed as a
single logical statement. Curly braces determine a ``scope'', inside
of which variable names may be re-used without ambiguity. A variable
defined between curly braces is said to be \emph{local} (to the scope
in which it is defined); its value cannot be used out of scope.

Function bodies are code blocks, as are the alternatives associated to
control statements. The compiler is not picky about spaces, tabs, and
newlines, so an input file should be organized in a way that makes the
file easy to read. Indentation signifies levels of nesting within code
blocks, but specific details are the focus of passionate debate. As
with variable naming, clarity and consistency are the important
criteria.


\subsection*{Headers and Pre-Processing}

A \CXX\ source file is compiled in multiple stages that occur
transparently to the user. The first step, pre-processing, involves
simple text replacement for file inclusion, macro expansion and
conditional compilation. Next, the source is compiled and assembled:
Human-readable language instructions are parsed, then represented in
assembly language. Finally, the object files are linked: Function
calls are resolved to hard-coded file offsets, possibly involving
external library files, and the program instructions are packaged into
an executable binary that the operating system can run.

Pre-processing is used much less in~\CXX\ than in~\code{C}; the
language itself supports safer and more featureful alternatives to
macros, such as \code{const} variables and inline functions. File
inclusion and conditional compilation are the chief uses of the
pre-processor. Lines of the form 
\begin{verbatim}
  #include <cstdlib>
  #include "epix.h"
\end{verbatim}
cause the contents of a \emph{header file} to be read into the source
file. A header file contains variable and function
\emph{declarations}, statements that specify types and names but do
not define actual data. Declarations tell the compiler just enough
to resolve expressions and function calls without knowing specific
values or function definitions. 

Conditional compilation is similar to conditional \LaTeX\ code, and is
best explained by example. A file might be used to produce both color
and monochrome output as follows:
\begin{verbatim}
//#define COLOR  // uncomment for color
#ifdef COLOR
  ...  // code for generating color figure
#else
  ...  // monochrome code
#endif // COLOR
\end{verbatim}
The ``compiler symbol'' \code{COLOR} is an ordinary \CXX\ name;
compilation is controlled by commenting or uncommenting the
\code{\#define} line. Multiple decisions on the same symbol may appear
in a file. An \code{\#else} block is optional, but every
\code{\#ifdef} must have a matching \code{\#endif}. Commenting the
\code{\#endif} is a good habit; in a realistic file, the start and end
of a conditional block may be separated by more than one screen.


\subsection*{Comparison with \LaTeX\ Syntax}

As a programming language, \CXX\ provides certain features common to
all languages (such as \LaTeX, Metapost, Perl, Lisp\ldots) and adheres
to rules of grammar. Salient differences between \LaTeX\ and \CXX\ 
include:

\begin{enumerate}
\item Every \CXX\ statement and function call must end with a
  semicolon. An omitted semicolon may result in a cryptic error
  message from the compiler. Pre-processor directives, which start
  with a~\code{\#}, do not end with a semicolon.

\item Backslash is an escape character in \CXX:\index{Labels}
\begin{verbatim}
  // Put label $y=\sin x$ at (2,1)
  // Note single  ^ backslash in output
  label(P(2,1), P(0,0), "$y=\\sin x$"); 
  //       Double backslash ^^ in source
\end{verbatim}
  
\item Variable and function names may contain letters (including
  underscore) and digits \emph{only}, are case sensitive, and must
  begin with a letter.

\item Variables in \CXX\ must have a declared \emph{type}, such as
  \code{int} (integer) or \code{double} (double-precision floating
  point). If a variable has global scope and its value does not
  change, the definition should probably come in the preamble or at
  the beginning of \code{main}. Local variables should be defined in
  the smallest possible scope. Unlike~\C, \CXX\ allows variables to be
  defined where they first appear.
  
\item \CXX~requires explicit use of~\code{*} to denote multiplication;
  juxtaposition is not enough. \CXX~does not support the use of
  \verb+^+ for exponentiation, e.g., \verb+t^2+ is invalid. Instead,
  use \code{t*t}~or \code{pow(t,2)}.

\item \CXX\ has single- and multi-line comments. Everything between a
  double slash and the next newline is ignored, while the strings
  \verb+/*+~and \verb+*/+ delimit multi-line comments. A single-line
  comment may appear within a multi-line comment, but the compiler
  does not nest multi-line comments.

\end{enumerate}

Between them, \C~and \CXX~have about 100~reserved keywords which
\textbf{cannot} be used as function or variable names. The script
\code{keywords} packaged with \ePiX\ is a simple lookup utility, meant
to help you avoid name clashes. To find keywords containing the string
\code{type}, for example, do ``\code{keywords type}''.
  

\section{The Camera}

\index{Camera|(emph}\index{Screen|(}\index{Camera!lens|(}
\ePiX\ depicts a Cartesian world by projecting mathematically to a
screen plane, then affinely scaling to a printed page. The camera,
which maps the world to the screen, consists of a \emph{body} (data
that determines the position and orientation of the camera) and a
\emph{lens} (the actual mapping to the screen plane).


\subsection*{The Body}

The camera's spatial orientation is described by a triple of mutually
perpendicular unit vectors.  In memory of happy days at the beach,
these vectors are called \emph{sea}, \emph{sky}, and~\emph{eye}. The
screen plane is parallel to the sea-sky plane; the sea vector points
horizontally to the right, sky points vertically upward. The eye is
their cross product, which points directly at the viewer.

The sea-sky-eye basis is located at the \emph{target}, so the target
is the origin of the screen plane. The \emph{viewpoint} lies on the
line through the target in the direction of the eye vector, and is the
center of projection for the default lens.  The distance from the
viewpoint to the target is the \emph{range}.  The orientation,
viewpoint, target, and range completely (and redundantly) determine
the camera's geometric situation in the world.


\subsection*{The Lens}

The \emph{lens} is a mapping from the world to the screen.  \ePiX\ 
comes with three lenses: \emph{shadow} (the default), \emph{fisheye},
and \emph{bubble}. Each lens simulates the appearance of world objects
as seen by an observer at the viewpoint. The shadow lens is point
projection from the viewpoint to the screen plane.  Each of the other
lenses performs radial projection to a sphere centered at the
viewpoint, then maps the sphere to the screen plane; the fisheye lens
does orthogonal projection (so the entire image lies inside a disk,
and positions behind the camera are inverted) while the bubble lens
does stereographic projection from the point directly behind the
viewpoint.
\index{Screen|)}

Other lenses may be defined; the file \filename{camera.cc} may be
consulted as a template. Syntactically, a lens is a \code{pair}-valued
function of a single \code{P}~argument. To ensure expected behavior a
lens should respect the meaning of the camera body.
\index{Camera!lens|)}


\subsection*{Manipulating the Camera}

\index{Camera!manipulation|(}
The camera is a \CXX\ class, manipulated by member functions.  The
\code{begin()} command initializes the camera, so all changes of
camera must occur in the body of the figure. By default, the viewpoint
is at large distance on the positive $x_3$-axis, looking down on the
$(x_1,x_2)$-plane. This provides the expected behavior for
2-dimensional figures.

The viewpoint and target may be set individually. In general, either
operation changes the direction of the eye vector, which forces \ePiX\ 
to re-determine the sky vector. When possible, the $x_3$-axis projects
parallel to the sky; otherwise the $x_2$-axis is used.  The viewpoint
and target can be moved along the eye axis, changing the range while
preserving the orientation: \code{range()} fixes the \code{target},
and \code{focus()} fixes the viewpoint. Each command re-sizes the
image; note that increasing the focus \emph{enlarges} the image.

The camera may be rotated about any of its axes. Axes of rotation pass
through the target, so rotations about the sea or sky vectors change
the viewpoint. For best control of the camera, set the target first,
then the viewpoint. If desired, perform eye rotations last.
\index{Camera!manipulation|)}
\index{Camera|)}


\section{Clipping and Cropping}

\index{Clipping|(}
\index{Cropping|(}
\ePiX\ provides two masking operations to handle figure elements that
lie far from the target: clipping (in the world) and cropping (in the
screen).  The ``clip box'' may be regarded as a set of ``walls''.
When clipping is active, objects outside the walls are not visible. By
default, the clip box is a large cube centered at the origin.

The ``crop box'' is a rectangle in the screen plane.\index{Screen}
When cropping is active, objects that project outside this rectangle
are not visible.  By default, the crop box is the bounding
box.\index{Bounding Box} Since the figure is drawn (on the page) by
affinely scaling the bounding box to a specified \LaTeX\ box, default
cropping ensures that a figure lies inside the printed region
allocated by \LaTeX.
\begin{figure}[hbt]
  \begin{center}
    \input{torus.eepic}\hspace*{0.75in}
    \input{torus_clip.eepic}\qquad
    \input{torus_crop.eepic}
  \end{center}
  \caption{Clipping and cropping a torus mesh (boxes added).}
  \label{fig:crop}
\end{figure}

By default, clipping and cropping are off. The command
\code{clip(bool)} (de)activates clipping. The argument defaults to
\code{true}. The clip box may be set with the commands
\begin{verbatim}
  clip_box(P pt1, P pt2); // opposite corners
  clip_box(P pt);         // opposite corners pt and -pt
  clip_to (P pt);         //                  pt and P(0,0,0)
\end{verbatim}
Analogously, cropping is (de)activted with \code{crop(bool)}. The crop
box is given by a pair of opposite corners, \code{crop\_box(pt1,pt2)};
third coordinates are discarded. With no arguments, the command
\code{crop\_box()} re-sets the crop box to the bounding box.
\index{Clipping|)}
\index{Cropping|)}


\section{Attributes}

At a minimum, \ePiX\ must be told how large the printed figure will
be, and how large a Cartesian rectangle to allocate. The preamble must
contain enough information to create a working state.  In the body of
an input file, the ``drawing state'' determines the figure's
appearance. Attributes are declarations, set by commands that accept
arguments of the stated type, possibly \code{void}. Color and path
width are controlled by writing \LaTeX\ commands to the output file
immediately; the remaining attributes are managed internally.

\begin{itemize}
\item Angular mode: \code{radians()}, \code{degrees()}, or
  \code{revolutions()}.

\item Path thickness: \code{plain()}, \code{bold()},
  \code{pen(double)}.\index{Path!width}

\item Path style:\index{Path!style}
  \begin{itemize}
  \item \code{solid()}
  \item \code{dashed(double)}. Optional argument is the dash density,
    the fraction (0.05--0.95) of the path taken up by dashes.
  \item \code{dotted(double)}. Optional argument is the diameter
    (in~\code{pt}) of a \code{dot}. 
  \end{itemize}
  The commands \code{dash\_length(double)} and \code{dot\_sep(double)}
  set the distance (in~\code{pt}) between vertices of a dashed or
  dotted path.
  
\item Color: \code{rgb(densities)}, \code{cmyk(densities)},
  \code{primary(density)}. 
%% A density is a \code{double} between~0 (no color) and~1 (full
%% saturation).

\item Filling: \code{fill(bool)}, argument defaults to \code{true}.
  \index{Filling}
  \begin{itemize}
    \item Gray depth: \code{gray(double)}, 0=white, 1=black.
    \item PSTricks (q.v.) filling style, fill color.
  \end{itemize}

\item Text rotation: \code{label\_angle(double)}\index{Labels!rotated}

\item Clipping and cropping (q.v.)

\item Camera (q.v.)
\end{itemize}


\subsection*{Angular Mode} 

\index{Angular mode}
\ePiX\ has three angular modes: \code{radians()} (the default),
\code{degrees()}, and \code{revolutions()}.  These modes affect all
trigonometric operations, including camera rotations, the drawing of
arcs and ellipses, polar plotting, label angle, and the trig functions
themselves. Angle-sensitive trig functions are capitalized, e.g.,
\code{Cos}, \code{Tan}.


\subsection*{Color and Shading} 

\index{Color|(} \ePiX\ provides color output via the \code{pstcol}
package, using the \code{rgb}~and \code{cmyk} models. Gray shading of
regions is supported through \filename{eepic.sty} (without requiring
\code{pstcol}).  Colors are best previewed by converting the document
to Postscript or PDF. Alternatively, EPS files can be previewed in
\code{xdvi}.

An \code{rgb}~color is determined by three floating-point densities
between~0 (no color) and~1 (full saturation).  A \code{cmyk} color
is similarly specified by four floats. Densities outside the
range~$[0,1]$ are ``clipped''.  Like line style, a color remains in
effect until superseded. Seven primary colors and white are available
by name. (Drawing in \code{white} can be used like correction fluid
to remove pieces of a figure accurately.)
\begin{verbatim}
  red();            // rgb(1,0,0);
  magenta(0.6);     // cmyk(0,0.6,0,0);
  rgb(0.2,0.7,0.8); // custom color
\end{verbatim}


\subsection*{PSTricks}

\index{PSTricks|(emph}\index{Filling!color|(}\index{Color!shading|(}
PSTricks is a powerful collection of macros, by Timothy
van~Zandt~\cite{Z} and others, for incorporating PostScript in \LaTeX.
\ePiX\ uses PSTricks primarily for color filling, but does not yet
work seamlessly. PSTricks should be used in a file only when
absolutely necessary. \textbf{The discussion below assumes PSTricks is
  needed in the file}.  The command
\begin{verbatim}
  use_pstricks(bool);
\end{verbatim}
sets an internal flag that determines whether a path is drawn as an
\filename{eepic} \code{path} or as a PSTricks
\code{psline}.\index{Path!in PSTricks} When issued before
\code{begin()}, this command also synchronizes the PSTricks
\code{unitlength} and line width with \ePiX, and sets the fill style
to solid.  PSTricks may be activated and deactivated freely, but
\emph{the first activation must occur in the preamble}.

While \code{pstricks} is active, commands of the form
\begin{verbatim}
  fill_color("<color name>");
  psset("<pstricks command>");
\end{verbatim}
are used to set attributes such as the style and color of paths and
filling. The default fill color is \code{"white"}.  This snippet,
taken from the sample file \filename{contour.xp}, shows how to define
a new color and use \code{psset()}.
\begin{verbatim}
  use_pstricks();      // synchronize length, etc.
  begin();
  use_pstricks(false); // temporarily disable
  ...
  std::cout << "\n\\newrgbcolor{orange}{1 0.7 0.2}";
  psset("fillcolor=orange, linecolor=green, linewidth=1.5pt");
\end{verbatim}
The PSTricks manual should be consulted for information.

Two major incompatibilities involve and filling and color. Unless the
fill style is set explicitly to \code{none}, PSTricks fills all paths,
even if they are not closed. Second, PSTricks manipulates colors by
named strings, so raw output is required to exploit the full power of
PSTricks from \ePiX.
\index{Color!shading|)}\index{Filling!color|)}\index{PSTricks|)}
\index{Color|)}


\section{The Path Class}

\index{Path|(emph}
A \code{path} data structure is \ePiX's low-level ordered list of
points that can be cropped, clipped, mapped, concatenated, and drawn.
Raw path data is useful for complicated paths built in pieces.
Available constructors are:
\begin{verbatim}
  path(p1, p2);          // line (endpoints)
  path(p1, p2, p3);      // quadratic spline
  path(p1, p2, p3, p4);  // cubic spline
  path(p1, v1, v2, t_min, t_max); // cf. ellipse arc
  path(f, t_min, t_max); // graph or parametrized path

  polyline(n, &p1, ...); // n points, followed by pointers
  polygon(n, &p1, ...);  // same, but marked as closed
\end{verbatim}
The first argument of the parametrized path constructor is a real- or
vector-valued function of one variable. In each of the first five
constructors, an optional final argument may be supplied to specify
the number of points used. 

The polyline and polygon constructors accept an unknown number of
arguments; consequently, their arguments must be passed as
``pointers'', or memory addresses:
\begin{verbatim}
  polyline(2, P(0,0), P(1,1)); // wrong; object arguments

  P p1 = P(0,0), p2=P(1,1);
  polyline(2, &p1, &p2);       // right; pointer arguments
\end{verbatim}
This makes polylines and polygons inconvenient for quick-and-dirty
use, but imposes little burden if the vertices have been defined
elewhere, as they might be in a logically structured file.

Compound paths may be built by concatenation. If \code{path1}~and
\code{path2} are paths, then the commands
\begin{verbatim}
  path1 += path2;
  path1 -= path2;
\end{verbatim}
replace \code{path1} with the result of traversing \code{path1}
``forward'', then following \code{path2} in the forward or reverse
direction (respectively). The notation is meant to suggest
1-dimensional homology chains. The sample file \filename{contour.xp}
illustrates path creation and manipulation.  Finally, note that a
\code{path} is a data structure, not a drawing command. The
\code{path::draw()} function must be called explicitly to create
visible output.


\subsection*{Path-Like Objects}

Path-like objects comprise polygons with a fixed number of vertices
(lines, triangles, quadrilaterals) and objects built from them
(including coordinate axes, grids, and arrows); curves with a variable
number of points (ellipses, arcs, splines); and plots of functions of
one variable.

Fixed-data polygons are drawn with high-level commands.
\begin{verbatim}
  line(p1, p2);
  triangle(p1, p2, p3);
  quad(p1, p2, p3, p4);
  rect(p1, p2);             // coordinate rectangle
  spline(p1, p2, p3);       // quad spline with control pts
  spline(p1, p2, p3, p4);   // cubic spline
\end{verbatim}
A \code{line()} accepts an optional numerical argument that acts as
an expansion parameter: \code{line(p1,p2,t);} draws a segment with
midpoint at the midpoint of \code{p1}~and \code{p2}, but having
length scaled by~$2^{t/100}$. (That is, $t=100$ doubles the length,
while $t=-100$ halves the length.) The arguments of \code{rect()}
must lie in a plane parallel to a coordinate plane.

\index{Filling|(}
Internally, \ePiX\ marks paths as closed or not; triangles,
quadrilaterals, polygons, ellipses, and arrows are closed. If a solid,
closed path is drawn while filling is active, the path is filled with
the current shade of gray. If filling is deactivated, the path is
drawn but not filled. Dashed and dotted paths cannot be filled in one
step; instead, fill the solid path, then draw the dashed/dotted
boundary.

The shade of gray is a number between~0 (white) and~1 (black), and
defaults to~0.3. Shading is opaque in PostScript, so the order of
figure elements is significant in the input file when filling is
active. 

\index{Filling!color}\index{PSTricks}
Color filling is available only through PSTricks. The features of
PSTricks that are relatively well-supported in \ePiX\ are illustrated
in the sample file \filename{contour.xp}. In principle, any PSTricks
features can be obtained through raw output. However, this approach is
generally discouraged since, for example, PSTricks color declarations
are not recognized by \eepic, and \emph{vice-versa}. The Next
Generation of \ePiX\ will work more seamlessly with PSTricks.
\index{Filling|)}

Elliptical arcs are specified by their center, a pair of vectors, an
angular range, and an optional number of points:
\begin{verbatim}
  ellipse(ctr, v1, v2, t_min, t_max, n);
\end{verbatim}
uses $n+1$ points to draw the parametrized path
\begin{equation*}
  t\mapsto \mathrm{ctr}+(\cos t)v_1 + (\sin t)v_2,\qquad 
  t_\mathrm{min}\leq t\leq t_\mathrm{max}
\end{equation*}
If omitted, the number of points defaults to~80 for a full turn, with
proportionally fewer points for an arc. When the angular range
subtends one or more full turns, the ellipse is marked as closed, and
will be filled if filling is active.

The circular arc parallel to the $(x_1,x_2)$-plane, having center~$p_1$
and radius~$r$, and subtending the angle (counterclockwise, in current
angle units) from $\theta_1$~to $\theta_2$ is drawn with
\begin{verbatim}
  arc(p1, r, theta1, theta2);
  arc_arrow(p1, r, theta1, theta2); // same, with arrowhead
\end{verbatim}
\index{Arrow}
If $\theta_2$ is smaller, the arc goes clockwise. The arrowhead goes
at~$\theta_2$. If an \code{arc\_arrow} is too short, nothing is drawn.
\index{Path|)}


\section{Geometric Data Structures}

Figure elements have an abstract description and a typographical
appearance. This section describes the former.  \ePiX\ provides
classes, \code{P}, \code{segment}, \code{circle}, \code{plane}, and
\code{sphere}, that can be used for Euclidean geometry constructions.
Each class is specified by a small amount of data (e.g., a center,
radius, and unit normal vector for a \code{circle}) and provides
constructors, intersection operators, affine transformations, and a
\code{draw()} function.

\ePiX\ also provides features for spherical and hyperbolic geometry,
especially the ability to draw lines in the half-plane and Poincar\'e
disk models of the hyperbolic plane, and to draw latitudes,
longitudes, and parametrized curves on a sphere.


\subsection*{Triples and Frames}

Points and displacements in space are represented as ordered triples.
The name of the type is~``\code{P}'', though ``\code{triple}'' may
be used for backward compatibility. Spherical and cylindrical/polar
constructors are provided. Standard (and non-standard) algebraic
operations---addition/subtraction, scalar multiplication; dot, cross,
and componentwise products, and orthogonalization---can be performed
on triples.  When forming symbolic expressions involving triples,
scalars must be collected together at left, triples at right. If
necessary, use parentheses to force a particular association.
\begin{verbatim}
  P pt(x,y,z);        // define pt = (x,y,z)
  double u=pt.x1();   // first coordinate of pt, etc.
  P(x,y);             // same as P(x, y, 0);
  polar(r, theta);
  cis(t);             // polar(1, t), aka Cos(t) + i*Sin(t)
  sph(r, theta, phi); // theta=longitude, phi=latitude
  P(a,b,c)|P(x,y,z);  // dot product, ax+by+cz
  P(a,b,c)&P(x,y,z);  // componentwise product, P(ax,by,cz)
  p*q;                // cross product, p x q
  J(p);               // quarter turn about the x3-axis
  p%q;                // orthogonalization, p (mod q)
\end{verbatim}
Explicitly, \code{p\%q}~is the unique vector \code{p+k*q} that is
perpendicular to~\code{q}.

A \code{frame} is a set of three mutually perpendicular unit vectors.
The standard \code{frame} is the set \code{E\_1}, \code{E\_2},
\code{E\_3}. The constructor turns three non-coplanar vectors into a
\code{frame}:
\begin{verbatim}
  frame(); // the standard frame
  frame(P v1, P v2, P v3); // orthonormalize (v1, v2, v3)
\end{verbatim}
The third vector of the new frame is positively proportional
to~\code{v3}, the second vector is positively proportional to
\code{v2\%v3}, and the first vector is the cross product. Thus, a
\code{frame} is right-handed, and does not depend on~\code{v1}.


\subsection*{Intersection}

The concept of \emph{genericity} is central to understanding
intersections of geometric data structures in \ePiX. For a working
definition, two objects that are disjoint, tangent, or coincident
intersect ``non-generically''. (Geometers will note that this differs
substantially from the usual definition.)  \ePiX's intersection
operators throw exceptions when the operands are non-generic. If a run
of \code{epix} terminates with an error message, check that you are
not trying to intersect badly-formed or situated objects.

In \ePiX, a \code{segment} is extended into a line for purposes of
intersecting.  Table~\ref{table:intersect} lists types of (generic!)
intersections in \ePiX. Intersection is commutative, so only the top
half of the table is shown. Not all intersections are defined.
\begin{table}[hbt]
\begin{center}
  \begin{tabular}{c|cccc}
    \code{*} & 
    \code{segment} & \code{circle} & \code{plane} & \code{sphere} \\
    \hline
    \code{segment} & \code{P} & \code{segment} & \code{P} & --- \\
    \code{circle}  &    & \code{segment} & \code{segment} & --- \\
    \code{plane}   &    &    &   \code{line} & \code{circle} \\
    \code{sphere}  &    &    &    & \code{circle}
  \end{tabular}
  \caption{\ePiX's intersection types.}
  \label{table:intersect}
\end{center}
\end{table}


\subsection*{Segments}

A \code{segment} is an \emph{unordered} pair of points. A segment may
be translated by a \code{P}, and the midpoint taken:
\begin{verbatim}
  segment L1 = segment(P(0,0), P(2,4));
  P mid = L1.midpoint();
  segment L2(mid, P(-2,3)); // form segment
  L2 += P(1,0);             // translate L2 by (1,0)
  L1.draw(); L2.draw();
  dot(L1*L2);               // point of intersection
\end{verbatim}


\subsection*{Circles}

A \code{circle} data structure consists of a center, radius, and a
perpendicular unit vector. Three constructors are provided:
\begin{verbatim}
  circle(P center, double radius, P normal);
  circle(P center, P point);
  circle(P p1, P p2, P p3);
\end{verbatim}
Unspecified trailing arguments in the first constructor take the
following defaults: If \code{normal} is not given, it is set
to~\code{E\_3}; if in addition, the radius is unspecified, it is
unity; finally, the center is the origin if no arguments are supplied.
As usual in~\CXX, only trailing arguments may be left implicit; the
call \code{circle(center, normal)} does \emph{not} create a unit circle
with the given center and normal.

The second creates a \code{circle} parallel to the $(x_1,x_2)$-plane
with the given center and passing through the given point. The third
returns the \code{circle} passing through the given points.
Exceptions are thrown if the \code{center}~and \code{point} do not lie
in a plane parallel to the $(x_1,x_2)$~plane, or if the three
points~$p_i$ are collinear.

A \code{circle} may be translated by a \code{P} with the
$+$~operator, or scaled by a \code{double} using the $*$~operator:
\begin{verbatim}
  circle C1=circle();      // unit circle
  circle C2 = C1+P(1,0.5); // translate up and right
  C2 *= 1.5;               // multiply radius by 1.5
  C1.draw(); C2.draw();    // view handiwork
  (C1*C2).draw();          // draw segment of intersection
\end{verbatim}


\subsection*{Spheres}

A \code{sphere} is specified by a point and a radius, by default the
origin and unity. Constructors are provided for a \code{sphere} of
specified center that passes through a given point, and for a
\code{sphere} given by a pair of antipodes:
\begin{verbatim}
  sphere(center, point);
  poles(north, south);
\end{verbatim}
As for \code{circle}s, translation and scaling operators are provided.
Capabilities specific to geography and spherical geometry are
described in Section~\ref{section:non-eucl}.
  
The \code{draw()} function of a \code{sphere} draws the horizon
visible from the current viewpoint. While this horizon is a circle in
object space, its image in the screen plane is generally an ellipse.
Further, antipodal points are not generally mapped to points that are
symmetrically placed with respect to the center of this ellipse. These
effects are most pronounced when the viewpoint is close to the
\code{sphere} and the center is not close to the \code{target}.

  
\subsection*{Lines and Planes}

\ePiX\ does not have a separate data structure for lines, but it does
provide a function, \code{Line}, that draws the line joining a pair of
points. Naturally, only a segment can be drawn. \ePiX\ always
\code{crop}s a \code{Line}, and in addition removes the half-line that
lies behind the observer. Thus, a \code{Line} appears as a printed
segment, either with ends on the bounding box, or with at most one end
in the interior, representing a point on a visible horizon. The Next
Generation (a.k.a.\ \code{ePiX3d}~or ``Version~2.x'') will supply a
line data structure.

A \code{plane} is specified by a point and a unit normal vector, or by
three non-collinear points. The \code{draw()} function renders the
lines of intersection of the \code{plane} with the faces of the
clip box. Unless the clip box has been set manually, these lines of
intersection will almost surely avoid the bounding box, and therefore
be invisible.


\section{Domains and Plotting}

\index{Plotting|(emph} To fix terminology, the noun ``map'' will be
used for a \CXX\ function that accepts one or more \code{double}
arguments or a \code{P} argument, and returns a \code{double}~or
a~\code{P}. Maps can be depicted in two mathematically useful ways:
graphing (which retains information about the domain), and drawing the
image as a parametrized curve or surface (which discards domain
information).  Usually, a \code{double}-valued map is graphed, while a
\code{P}-valued map is viewed parametrically. The verb ``plot'' is
used generically to mean either sort of depiction, and the noun
``plot'' refers to the result of plotting.

\index{Domain|(emph}
The ``domain'' of a map is the set of allowed input values. In \ePiX,
a \code{domain} is a coordinate box of dimension one, two, or three.
For uniformity, a 2-dimensional domain is actually a 3-dimensional
domain whose ``thickness'' along the $x_3$-axis is~0. Similarly, a
1-dimensional domain is a ``highly degenerate'' 3-dimensional domain.
Domains facilitate the plotting of families of mappings, and the
selective plotting of parts of a mapping; \code{domain} operators for
these purposes are described below. 

A plot depends on a map and a domain. For example, suppose~\code{f} is
a \code{P}-valued map of two (\code{double}) variables, and~\code{R}
is a two-dimensional domain. To construct a plot, subdivide~\code{R}
with a ``coarse'' rectangular mesh, say having $m_1\times m_2$
sub-rectangles.\index{Domain!coarse mesh|(} For each ``grid line'', plot
the parametric curve obtained by evaluating~\code{f} along the line.
The result is the (wiremesh) plot of~\code{f} over~\code{R}.

A detail was omitted in the preceding paragraph: the number of points
plotted in each parametric curve. Many graphing programs use as many
points as dictated by the coarse mesh, so that the surface is made of
quadrilaterals (Figure~\ref{fig:plot}, left). In \ePiX, a domain
possesses a ``fine'' mesh, say that divides~\code{R} into $n_1\times
n_2$ subrectangles (Figure~\ref{fig:plot}, right).
\index{Domain!fine mesh|(} The individual parametric curves comprising
the surface are drawn using the fine subdivisions. A plot can conform
to the ``true'' surface without requiring many quadrilaterals.

\begin{figure}[hbt]
  \begin{center}
    \input{plot_a.eepic}\hspace*{0.5in}
    \input{plot_b.eepic}
    \caption{A surface with coarse ($6\times18$) and fine
      ($12\times60$) meshes.}
    \label{fig:plot}
  \end{center}\index{Plotting!wiremesh}
\end{figure}


A \code{domain} is defined by giving a pair of opposite corners, the
coarse mesh, and the fine mesh. A domain's meshes are independent; the
fine mesh need not be a ``multiple'' of the coarse mesh.  Indeed, the
``fine'' mesh may be \emph{more coarse} than the ``coarse'' mesh.
Usually, however, the fine mesh is a ``small multiple'' of the coarse
mesh.
\index{Domain!fine mesh|)}
\index{Domain!coarse mesh|)}

\index{Domain!resizing}
A domain can be resized in any coordinate direction for which the
thickness is positive, using object-oriented syntax. For the
domain~\code{R} above,
\begin{verbatim}
  domain R_new = R.resize2(0.5,0.75);
\end{verbatim}
defines a new domain having opposite corners $(0,0.5)$~and $(1,0.75)$.
To the extent possible, resizing attempts to preserve absolute grid
sizes. In this example, the new domain has $6\times4$~coarse mesh
($18/4=4.5\to4$) and $12\times15$~fine mesh.  It's best to choose
meshes so that resizing does not cause integer truncation.

\index{Domain!slicing}
A domain can be sliced by setting one variable to a constant; the
result is a domain whose dimension is one smaller than the original:
\begin{verbatim}
  domain R1 = R.slice1(0.3);
\end{verbatim}
creates the 1-dimensional domain with endpoints $(0.3,0)$~and
$(0.3,1)$, and having coarse mesh~18~and fine mesh~60.  Resizing and
slicing commands may be used directly in a \code{plot} command:
\begin{verbatim}
  plot(f, R.slice1(0.3));
\end{verbatim}
draws the parametric curve obtained by slicing~\code{R} at $x_1=0.3$,
then applying~\code{f}. 

For plotting families of maps, a \code{domain} has ``\code{slices}''
operators that return the list of \code{domain} slices obtained by
setting one variable to components of the coarse mesh:
\begin{center}
  \input{slices.eepic}
\end{center}

When plotting a function of one variable, constructing a \code{domain}
is possible but tedious.  The command
\begin{verbatim}
  plot(f, t_min, t_max, n);
\end{verbatim}
plots~\code{f} over the interval \code{[t\_min, t\_max]} using
\code{n}~subintervals (\code{n+1}~points).

The ability to slice or resize makes \code{domain}s more useful when
plotting functions of two or three variables.  The commands below
typify plot commands; $f$~is a function of one variable, and~$F$ is a
function of 2~or 3~variables. If the function is real-valued, the
graph is drawn; otherwise the image is drawn.
\begin{verbatim}
  // two ways to plot F on [a,b] x [c,d]
  plot(F,  P(a,c), P(b,d), mesh(N1,N2), mesh(n1,n2));

  domain R(P(a,c), P(b,d), mesh(N1,N2), mesh(n1,n2));
  plot(F, R);

  plot(F, R.slice1(a));      // plot F on part of the boundary
  plot(F, R.slices1());      // plot F over all slices x1=const
\end{verbatim}
A function of 3~variables can be plotted over a 1-~or 2-dimensional
domain, while a function of 2~variables cannot be plotted over a
3-dimensional domain.
\index{Domain|)}
\index{Plotting|)}


\subsection*{Utility Functions}

In this section, \code{f}~and \code{g} are \code{double}-valued
functions of one variable. \ePiX\ defines numerical functions that
return the maximum or minimum value on an interval, approximate the
location of roots, and perform calculations with derivatives and
definite integrals.
\begin{verbatim}
  sup(f, a, b);     // max/min of f on [a,b]
  inf(f, a, b);
  newton(f, g, x0); // find approximate crossing point
\end{verbatim}
Newton's method returns the crossing point of the given functions,
starting from the specified seed. If a critical point is hit or
20~iterations pass, a warning is issued and the current result
(probably incorrect) is returned. The second function~$g$ defaults to
the zero function if omitted. 

The classes \code{D}~and \code{I} are used to calculate values of
derivatives and integrals, and to plot these functions.
\begin{verbatim}
  D df=D(f);   // data structure representing the derivative
  df.eval(t);  // return f'(t)
  df.left(t);  // deriv from left at t:  (f(t)-f(t-dt))/dt
  df.right(t); // deriv from right at t: (f(t+dt)-f(t))/dt

  I prim=I(f,a); // representation of the integral from a
  prim.eval(b);  // numerical integral of f over [a,b]
  double val=I(f).eval(1);   // \int_0^1 f
\end{verbatim}
The lower limit on an integral is~0 by default.

Tangent lines and envelopes (families of tangent lines) are drawn with
\begin{verbatim}
  tan_line(f, t);                // f real- or vector-valued
  envelope(f, t_min, t_max, n);  // family of tangent lines
  tan_field(f1, f2, t_min, t_max, n); // field of tangents
\end{verbatim}
The sample files \filename{conic.xp}~and \filename{lissajous.xp}
illustrate these features.


\subsection*{Calculus Plotting} 

\index{Plotting!calculus|(}
\index{Plotting!vector fields|(}
\ePiX\ can plot the derivative or definite integral of real-valued
functions, solve ODEs in two or three variables, and graph slope- or
vector fields. Let \code{f}~be a real-valued function of one variable,
\code{F}~a \code{P}-valued function of two or three variables.
\begin{verbatim}
  plot(D(f), a, b, n);  // plot f' over [a,b]
  plot(I(f, x0), a, b, n);

  ode_plot(F, p_0, t_min, t_max, n); 
  flow(F, p_0, t_max, n);
\end{verbatim}
The second command graphs the definite integral $x\mapsto\int_{x_0}^x
f(t)\,dt$ over~$[a,b]$. As above, $x_0$~defaults to~0.

The third command plots the solution curve of the initial-value
problem $\dot{x}=F(x)$, $x(0)=p_0$, over the specified time interval.
If $t_\mathrm{min}$ is omitted, its value is~0, so the curve starts
at~$p_0$. With manual calculation to rotate a planar field a quarter
turn, \code{ode\_plot} can be used to draw level curves (isobars) of a
function; see the sample file \filename{dipole.xp}. The fourth command
returns the result of starting at~$p_0$ and flowing by~$F$ for
time~$t_\mathrm{max}$, using Euler's method with~$n$ timesteps.  This
is useful for placing markers or arrowheads precisely along a flow
line. A vector field itself may be drawn in three ways:
\begin{verbatim}
  vector_field(F, p, q, n1, n2); // proportional length
  dart_field  (F, p, q, n1, n2); // constant length
  slope_field (F, p, q, n1, n2); // directionless, const length
\end{verbatim}
The field is sampled at the grid points of the coordinate rectangle
whose corners are specified.
\index{Plotting!vector fields|)}
\index{Plotting!calculus|)}


\subsection*{Recursive Fractal Curves}  

\index{Path!fractal|(}
Consider a path made up of equal-length segments that can point at any
angle of the form~$2\pi k/n$ radians, for $0\leq k<n$, like spokes on
a wheel. A path is specified by a finite sequence of integers, taken
modulo~$n$. For example, if $n=6$, then the sequence $0, 1, -1, 0$
corresponds to the ASCII path \verb+_/\_+.  \ePiX's fractal
approximation starts with such a ``seed'' then recursively (up to a
specified depth) replaces each segment with a scaled and rotated copy
of the seed. The seed above generates the standard von~Koch snowflake
fractal. In code:
\begin{verbatim}
  const int seed[] = {6, 4, 0, 1, -1, 0};
  fractal(P(a,b), P(c,d), depth, seed);
\end{verbatim}
The first entry of \code{seed[]} (here~6) is the number of
``spokes''~$n$, the second~(4) is the number of terms in the seed, and
the remaining entries are the seed proper. The final path joins
$(a,b)$~to $(c,d)$. The number of segments in the final path grows
exponentially in the depth, so depths larger than 5~or 6 are
likely to exceed the capabilities of \LaTeX\ and/or PostScript.
\begin{figure}[hbt]
  \begin{center}
  \input{koch.eepic}
  \end{center}
  \caption{Successive iterations of \{\code{4,8,0,1,0,3,3,0,1,0}\}}
\end{figure}
\index{Path!fractal|)}


\section{Non-Euclidean Geometry}
\label{section:non-eucl}

Hyperbolic line segments are specified by their endpoints in the upper
half space or ball (Poincar\'e) models.  In each case there is no
output if either endpoint lies outside the model.
\begin{verbatim}
  hyperbolic_line(p, q);
  disk_line(p, q);
\end{verbatim}
For compatibility with 2-dimensional hyperbolic space, the half-space
model is the set $\{(x_1,x_2,x_3)\mid x_2>0\}$.

\index{Plotting!spherical|(}
A \code{frame} determines geographical coordinates on a \code{sphere}:
the first element points toward longitude~0 on the equator, the third
element points to the north pole.  A latitude line depends on a
\code{sphere}, a \code{frame}, the numerical latitude, and a range of
longitudes. A longitude line is described similarly.
\begin{verbatim}
  latitude(lat, long_min,long_max, sphere S, frame coords);
  longitude(lngtd, lat_min, lat_max, sphere S, frame coords);
\end{verbatim}
By default, \code{coords} is the standard \code{frame} and \code{S} is
the unit \code{sphere}. These commands draw only the portion of the
curve that is visible from the current \code{viewpoint}. The function
\code{back\_latitude} draws the invisible portion of a latitude line.

Parametrized paths on the unit sphere can be specified either by
radial projection of a space curve, or by stereographic projection of
a plane curve:\index{Path!on sphere}
\begin{verbatim}
    plot_R(phi, t_min, t_max, n);    // radial
    plot_N(f1, f2, t_min, t_max, n); // from north pole
    plot_S(f1, f2, t_min, t_max, n); // from south pole
\end{verbatim}
Attempts to perform radial projection on a path through the origin
will generate division-by-zero errors. Stereographic projection maps
the equatorial plane $\{x_3=0\}$ to the unit sphere by projection from
the corresponding pole: $N=(0,0,1)$, $S=(0,0,-1)$.

Each spherical plot command accepts a prefix \code{front} or
\code{back} that prints only the portion of the path visible or
invisible (respectively) from the current \code{viewpoint}.

Because of the way \ePiX\ layers output, it is generally best to put
hidden portions of the input before visible portions, with line width
and/or style that suggests hidden lines.
\index{Plotting!spherical|)}


\section{Animation}

\index{Animation|(}
\ePiX\ is ideally suited to the creation of mathematically accurate
animations: If a figure depends suitably upon a ``time'' parameter,
then a loop can be used to draw the entire figure for multiple time
values, yielding successive ``snapshots'' of the figure as time
progresses. The shell script~\code{flix} automates the process of
compiling a suitable input file into a collection of \filename{png}s
and assembling these frames into a \filename{mng} animation.
ImageMagick is the image-handling engine.

A \code{flix} file is an \code{epix} file with two restrictions:
\begin{itemize}
  \item The \code{double} variable \code{tix} is used as
  ``clock''.

  \item \code{main} accepts two commandline arguments and sets
  \code{tix} accordingly.
\end{itemize}
Jay Belanger's \code{emacs} mode recognizes the file
extension~\filename{.flx} and inserts template code if an empty
buffer is opened. Creation of \code{flix} files is as easy as creation
of \code{epix} files.  The directory \filename{samples/extras}
contains a handful of \code{flix} files that may be consulted for
ideas.

A ``typical'' \code{.flx} file may take 30~seconds to a couple of
minutes to compile. To present the impression that work is being done,
\code{flix} prints a progress bar, counting the number of
\filename{eps} files that have been created. There will be a delay of
a few seconds after the last frame is produced, during which
ImageMagick's \code{convert} utility creates \code{png} files from
\filename{eps} files, then assembles the movie.

To facilitate debugging, \code{elaps} can be run on a \code{flix}
file. \code{elaps} runs in a fraction of the time, and if \code{elaps}
can't produce a viewable image, \code{flix} will surely fail.

By default, \code{flix} creates movies with 24~frames, in which
\code{tix} runs from~0 to~1, and animates at $0.08$~sec/frame.
Command-line options change these and other parameters; please use
\code{flix}'s built-in help for details.  \index{Animation|)}


\section{Troubleshooting}

Seven files comprise \ePiX: a header (\filename{epix.h}), compiled
library (\filename{libepix.a}), and five shell scripts (\code{epix},
\code{elaps}, \code{laps}, \code{flix}, and \code{keywords}). Four of
these files are generated at compile time, using variables found in
the \filename{Makefile}.  In addition to the executable portions of
the program, there are a few dozen sample files, configuration files,
\filename{README}s, and miscellaneous documentation.  These components
are placed into standard locations where they can find each other.

\subsection*{Installation Problems}
\index{Installation|(}

Public installation (in \code{/usr/local}) is the default. Unpack
the tar file, then \code{cd} to the source directory
(\filename{epix-x.y.z\_complete}). If you are installing in
\filename{/usr/local}, your \CXX~compiler is~\code{g++}, and
\code{bash} is in \filename{/bin/bash}, then do
\begin{verbatim}
make
make contrib (optional)
make test    (optional)
make install
\end{verbatim}
(the last as \code{root}, if necessary). If \code{bash} is not in
\filename{/bin/bash}, edit the script \filename{newbash.sh} and run it
in the source directory.  If the install directory or compiler are not
as above, edit the \code{Makefile} accordingly, then build. If a
failed partial compile creates the header file \filename{epix.h}, you
\emph{must}\/ do \code{make clean} before re-compiling.

If you have installed in \filename{\$HOME}, please see the
\filename{POST-INSTALL} file for additional instructions. 

The directory in which you install \ePiX\ is denoted
\filename{\$INSTALL}. A successful \code{make install} creates the
following files:
\begin{verbatim}
  $INSTALL/bin/{epix,elaps,laps,flix,keywords}
  $INSTALL/include/epix.h
  $INSTALL/lib/libepix.a
  $INSTALL/man/man1/epix.1
\end{verbatim}
Sample files, notes and README files, configuration snippets, and (if
the \code{complete} package is installed) this tutorial are
installed in subdirectories of \filename{\$INSTALL/share/epix}.


\subsection*{Known Issues}

\index{Mac OS X}
If you install \ePiX\ on MacOS~X with the December, 2002 Apple
Developer Tools, you must manually run \code{ranlib} on the compiled
library:
\begin{verbatim}
  sudo ranlib /usr/local/lib/libepix.a
\end{verbatim}

Buggy versions of the conversion script \code{ps2epsi} have been
distributed with recent versions of RedHat and Mandrake. The author
has seen two \emph{different} one-byte errors that cause \code{elaps}
to fail with a \code{sed} error. If \code{elaps} generates a
\code{sed} error, you have a buggy \code{ps2epsi}.  Please google for
the specific error message or, if you're good with regular
expressions, just fix the script yourself. (In both prior instances,
the problem was an unescaped special character.)


\subsection*{Runtime Errors}

Typos frequently cause errors when \code{epix} is run; the compiler
generally issues a helpful message, naming the troublesome line of the
input file and the type of error.

Two common sources of error in a syntactically correct input file are
camera mis-placement and non-generic intersections. If the camera is
too close to objects in the scene, the lens may try to divide by zero
(or by a small epsilon), resulting in \code{nan} errors (or very large
coordinates) in the output file. In this situation, \code{epix} will
succeed, but \code{elaps} will hang. By contrast, an error message of
the form
\begin{verbatim}
  /usr/local/bin/epix: line 275: 27333 Aborted ...
\end{verbatim}
signifies an uncaught exception, hence a non-generic intersection.

Though the shell scripts are extensively tested, they are a possible
source of runtime problems. (For example, \code{elaps} tries to pass
options to the compiler while reserving other options for its own use,
all while trying to deal intelligently with malformed commands.) If a
script does something unexpected, it's probably a Feature, not a Bug.
In any case, please notify the author of anomalous behavior.

\paragraph{Hanging scripts}
By default, the conversion scripts run silently. Because \LaTeX~does
not use separate channels to return output and error messages,
\code{laps} can hang if there is a problem with the \LaTeX\ file. If
this happens, try typing ``\code{s}'' to put \LaTeX\ into scroll mode.
If this fails to return the prompt after a few seconds, type
\verb+<ctrl>-C+ to kill \code{laps}, and re-run in verbose mode:
\begin{verbatim}
laps -vv <filename>
\end{verbatim}
This will cause \LaTeX's error messages to be printed normally.

\paragraph{Command not found} In order to use \ePiX, the directory
\filename{\$INSTALL/bin} must be in your \code{PATH} (see 
\filename{POST-INSTALL}); type\\
\code{echo \$PATH}\\
\noindent to see your \code{PATH}. If the directory
\filename{\$INSTALL/bin/} is not in your \code{PATH}, please read
\filename{POST-INSTALL} or ask someone knowledgeable at your site for
help; the procedure varies depending on what shell you use. In any
case, you will need to modify your shell's configuration file.

\paragraph{Permission denied} This is  unlikely, but conceivable. For
each component of the program, do a long listing, e.g.\ (with the
appropriate install directory)
\begin{verbatim}
  ls -l /usr/local/bin/epix
\end{verbatim}
The header and library must be readable, and the shell scripts and
directories must be readable and executable. From the install
directory, do
\begin{verbatim}
chmod 0755  bin  include  lib  bin/{epix,elaps,laps,flix,keywords}
chmod 0644  include/epix.h  lib/libepix.a  man/man1/epix.1
\end{verbatim}
\index{Installation|)}

If you still cannot get \ePiX\ to run, please send email to the
author. Generally, it is helpful to specify the operating system
(e.g., RedHat~8.2, or Debian~Potato on a G4~Powerbook), the version of
the \C~compiler (e.g.~\code{gcc-3.2}), the version of \ePiX, and
where you downloaded the sources (CTAN, the project home page, etc.).
Specific commands you typed and error messages may also be helpful. If
you don't know what an error message means, feel free to paste or
attach it to your email as plain text. The author does not have access
to an enormous variety of systems, and sometimes reported errors never
get explained, but to date no error has failed to be worked around.


\subsection*{\LaTeX\ Errors}

A few things can cause \LaTeX\ to stop with an error message when
reading an \eepic\ file written by \ePiX. The most common is the
appearance of \code{nan} (not a number) where \LaTeX\ expects a
number. This generally indicates division by zero or bad
exponentiation.

When a number is very small, \ePiX\ may write it in exponential
notation. If this happens, \LaTeX\ will pause with an error message
when it tries to read, e.g., \code{1.4142135e-14}. This bug has been
addressed; please send the author a bug report if you encounter this
behavior in \ePiX\ code. You can manually edit the \eepic\ file,
replacing underflows with~$0$. In this eventuality, it's wise to
rename the edited file, lest \ePiX\ overwrite your changes the next
time you run it.

Overflow errors in \LaTeX\ are possible if a point has coordinates
larger than~$2^{16}$; make sure you're not trying to plot the graph of
a pole or something similar.


\chapter{Advanced Topics}
\label{chapter:adv}

This chapter covers \emph{ad hoc} tricks and open-ended techniques
that require relatively more programming sophistication. You will
almost surely need an external \CXX\ reference if you do not speak the
language.


\section{Hidden Object Removal}

\index{Hidden object removal|(}
\index{Plotting!surface|(}
\ePiX\ writes the output file in the same order that objects appear in
the input. The order is significant because PostScript builds a figure
in layers: Objects are drawn over objects that come earlier in the
file. Shaded polygons can be used to obtain surprisingly effective
hidden object removal in surface meshes. The techniques are still
provisional, however; this section describes a method that works for
the author, though is not sufficiently refined for inclusion in the
source code.

The basic idea is to create a shaded quadrilateral class that knows
its distance to the camera.\index{Camera} To draw parametrized
surfaces, quadrilaterals are written in any convenient order to a
\CXX\ vector, then sorted by distance to the camera and printed to the
output file in decreasing order of distance. The gray density of a
mesh element depends on the cosine of the angle between the normal
vector and the vector from the camera to the element. The method works
acceptably well in practice: Several surfaces can be drawn, and the
effect is fairly realistic. If the surface meshes are fine enough,
intersections are accurate and not jagged. However, the output file
tends to be enormous, and is not effectively human-readable.

The quadrilateral class might look like this:

\begin{footnotesize}
\begin{verbatim}
class mesh_quad
{
private:
  P pt1, pt2, pt3, pt4;
  double distance;

public:
  mesh_quad() 
  {
    pt1=pt2=pt3=pt4=P();
    distance=camera.get_range();
  }

  mesh_quad(P f(double u, double v), double u0, double v0)
  {
    pt1=f(u0+EPS,v0+EPS);       // EPS = "shrinkage"
    pt2=f(u0+du-EPS,v0+EPS);
    pt3=f(u0+du-EPS,v0+dv-EPS);
    pt4=f(u0+EPS,v0+dv-EPS);
    
    P center = 0.25*(pt1 + pt2 + pt3 + pt4);
    distance = norm(center-camera.get_viewpt());
  }

  double how_far() const { return distance; }

  void draw() 
  { 
    P normal = (pt2 - pt1)*(pt4 - pt1);
    normal *= 1/norm(normal);

    double dens  = 0.75*(1-pow(normal|LIGHT, 2)/(LIGHT|LIGHT));
    gray(dens); 
    quad(pt1, pt2, pt3, pt4); 
  }
};
\end{verbatim}
\end{footnotesize}
To utilize \CXX's sorting algorithm, a class is defined to measure the
distance to a mesh element:
\begin{footnotesize}
\begin{verbatim}
class by_distance {
public:
  bool operator() (const mesh_quad& arg1, const mesh_quad& arg2)
  { return arg1.how_far() > arg2.how_far(); }
};
\end{verbatim}
\end{footnotesize}
Assuming \code{f1}~and \code{f2} are parametric surface maps, the
class above can be used as follows in the body of the figure:
\begin{footnotesize}
\begin{verbatim}
  std::vector<mesh_quad> mesh(0); // list of mesh_quads

  for (int i=0; i<N1; ++i)
    for (int j=0; j<N2; ++j)
      {
        mesh.push_back(mesh_quad(f1, -1+du*i, -1+dv*j));
        mesh.push_back(mesh_quad(f2, -1+du*i, -1+dv*j));
      }

  sort(mesh.begin(), mesh.end(), by_distance());

  for (unsigned int i=0; i<mesh.size(); ++i)
    mesh.at(i).draw();
\end{verbatim}
\end{footnotesize}
For a complete file, please see \filename{samples/extras/tori.xp}.
\index{Plotting!surface|)}\index{Hidden object removal|)}


\section{Extensions}

Thanks to a suggestion of Andrew Sterian, \ePiX\ is extensible. The
preferred method is to create external modules for use at run time,
analogous to a \LaTeX\ macro/style file. It is also possible to modify
the \ePiX\ source code itself before compiling; this is inflexible
(and requires administrator privileges), but can be useful for
changing system-wide defaults.  User extensions span a spectrum, from
header files that require only basic knowledge of~\CXX\ to separately
compiled libraries that substantially extend the capabilities of
\ePiX.


\subsection*{Header Files}

A \CXX~header file conventionally has suffix~\filename{.h}, as in
\filename{myheader.h}. To use this custom header, put a line
\code{\#include "myheader.h"} in your source file.

User definitions can be easily and robustly implemented with ``inline
functions''. Inline functions are superficially similar to macros, but
are far more safe and featureful (since they are handled by the
compiler rather than by the pre-processor).  Examples are
\begin{verbatim}
  inline void Bold(void) { pen(1.5); }
  inline void purple(void) { rgb(0.5, 0, 0.7); }
  inline void draw_square(double s) { rect(P(-s,-s),P(s,s)); }
  inline double cube(double x) { return x*x*x; }
\end{verbatim}
The keyword \code{void} signifies a function that does not return a
value, or (when used as a parameter) a function that does not accept
arguments. Inline function definitions are syntactically exactly like
ordinary function definitions, but \emph{must} occur in a header file
or in the source file where they are used. The examples above might be
used in an input file as follows:
\begin{verbatim}
  Bold();  
  draw_square(cube(1.25));
\end{verbatim}



\subsection*{Compiling}

The next few sections outline the creation of a ``static library'' on
GNU/Linux, and explain how to incorporate custom features at runtime.
For more details, the source code, \filename{Makefile}, and your
system's documentation should provide a good start.

A small library is usually written as a \emph{header} file, which
contains function declarations (also called ``prototypes''), and a
\emph{source} file, which contains the actual code. Conventionally
(under *nix), these files have extension \filename{.h}~and \filename{.cc}
respectively. Header and source files may ``include'' other header
files, to incorporate additional functionality.
\begin{verbatim}
/* my_code.h */
#ifndef MY_CODE
#define MY_CODE
#include <cmath>   // standard library math header
#include "epix.h"  // ePiX header
using ePiX::P;

namespace Mine {   // to avoid name conflicts
  // functions for special relativity
  double lorentz_norm(P);
  bool   spacelike(P);
} // end of namespace
#endif // MY_CODE
\end{verbatim}
This file exhibits two ``safety features''. The three \code{MY\_CODE}
lines prevent the file from being included multiple times. In a file
of this size, inclusion protection is overkill, but as your code base
grows and the number of header files increases, this protection is
essential. Second, the header introduces a ``Mine'' namespace. Inside
this namespace, two functions are declared as prototypes, giving the
function's return type, name, and argument type(s). A header file
should be commented fairly liberally, so that a year or two from now
you'll be able to decipher the file's contents.  For a longer file,
version and contact information, an overall comment describing the
file's features, and license information are appropriate.

Next, the corresponding source file; definitions are also placed into
the namespace, and must match their prototypes from the header file
exactly.
\begin{verbatim}
/* my_code.cc */
#include "my_code.h"
using namespace ePiX;

namespace Mine {
  double lorentz_norm(P arg)
  {
    double x=arg.x1(), y=arg.x2(), z=arg.x3(); // extract coords
    return -x*x + y*y + z*z;
  }
  bool spacelike(P arg)
  {
    return (lorentz_norm(arg) > 0); // true if inequality is
  }
} // end of namespace
\end{verbatim}
Copies of these files are included with the source code so you can
experiment with them. Next, the source file must be ``compiled'',
``archived'', and ``indexed''. In the commands below, the percent sign
is the prompt.
\begin{verbatim}
% g++ -c my_code.cc
% ar -ru libcustom.a my_code.o
% ranlib libcustom.a
\end{verbatim}
Please see your system documentation for details on command options
and what each step does. For linking (below), the name of the library
file must begin ``lib'' and have the extension~\filename{.a}. Once these
steps are successfully completed, put the library \filename{libcustom.a}
and header file \filename{my\_code.h} in your project directory. You're
ready to use the code in an \ePiX\ figure.


\subsection*{Runtime Linking}

The script~\epix\ allows input files to be linked with external
libraries at run time, when the input file is compiled into a
temporary executable.

\epix\ recognizes command line options and passes them verbatim to the
compiler. The most commonly used options are those of the form
\begin{verbatim}
   -I<include>     -L<libdir>     -l<lib>
\end{verbatim}
There must be no space between the option flags \code{-I},
\code{-L}, and~\code{-l} and their arguments.  For example, to
link \filename{figure.xp} against \filename{mylibs/libcustom.a}, run
the command
\begin{verbatim}
   epix -Lmylibs -lcustom figure
\end{verbatim}
The options \code{-I. -L.} tell the compiler to look in the current
directory for header and library files.  Compiler options may appear
in any order, but must come before the name of the input file; options
that come after the input file are silently discarded.

Compiler options may be placed in the configuration file
\filename{\$HOME/.epixrc}, with syntax as above. A line in the config
file that contains a pound sign~(\code{\#}) is a comment, no matter
where in the line the~\code{\#} appears. If any non-comment line
fails to start with a dash, the rest of the file is silently
discarded.  Command-line options are read before the config file.

% \clearpage
% \setcounter{chapter}{0}
% \renewcommand{\thechapter}{\Alph{chapter}}

\appendix
\chapter{Software Freedom}
\index{Free software|(} 

Academics in general, and mathematicians in particular, depend on Free
software in their work. A good case can be made that proprietary
software is contrary to the academic ethic. Issues of access aside, if
one does not know what exactly went into a program, then one cannot
fully trust the results that come out, any more than one can trust
(for purposes of scientific publication) results of a commercial
testing lab.  Access to the source code is not all that is required,
though. To promote the dissemination of information, users should be
granted the four freedoms laid out in the GNU General Public License:
\smallskip

\code{Free0}: To run a program for any purpose

\code{Free1}: To study how the program works, and adapt it to your needs

\code{Free2}: To redistribute copies of the program

\code{Free3}: To improve the program, and release improvements to the
public
%\smallskip

Just as theorems are not restrictively licensed, I believe that
software we use in our academic work should be licensed in a way that
encourages openness and sharing. Releasing software under a standard
commercial license agreement is (to me) the equivalent of publishing
the statement of a theorem, while keeping the proof secret, and
charging people for each citation of the theorem. Releasing source
code alone, without giving users the freedom to modify it for their
own needs, is analogous to publishing a proof, but forbidding
readers from using the ideas of the proof in their own work.

The ultimate purpose of software is to allow us to be productive and
creative. I hope that this modest program is, in conjunction with the
much larger efforts of others (especially Donald Knuth, Richard
Stallman, and the many people who have contributed to the authorship
of \LaTeX\ and its packages), useful to you in your mathematical work.

Please visit the Free Software Foundation, at
\verb+http://www.fsf.org+, to learn more about Free Software and how
you can contribute to its development and adoption.
\index{Free software|)}


\chapter{Acknowledgements}

\ePiX\ is built on the work of many people (unfortunately, most of
whom I am unaware). The following people have contributed, sometimes
unknowingly but always generously: 

\paragraph{Infrastructure} Donald E. Knuth, Conrad Kwok, Leslie
Lamport, Tim Morgan, Piet van Oostrum, Sunil Podar, Richard Stallman

\paragraph{Enhancements} Jay Belanger, Robin Blume-Kohout, Guido
Gonzato, Svend Daug{\aa}rd Pedersen, Andrew Sterian

\paragraph{Debugging, advice, and other assistance} Jay Belanger,
Felipe Paulo Guazzi Bergo, Robin Blume-Kohout, Patrick Cousot, Stephen
Gibson, Dov Grobgeld, Bob Grover, Jim Hefferon, Jacques L'helgoual,
Yvon Henel, Hartmut Henkel, Herng-Jeng Jou, Walter Kehowski, Ross
Moore, Thorsten Riess, Neel Smith, Michael Somos, Andrew Sterian,
Ryszard Tanas, Kai Trukenmueller, Torbjorn Vik, Wenguang Wang, Gabe
Weaver, Mariusz Wodzicki

\begin{thebibliography}{9}
%
\bibitem{KnR} B.~Kernighan and D.~Ritchie, \emph{The \code{C}
    Programming Language}, Second Ed., Prentice-Hall Software Series,
    1988
%
\bibitem{GNUC} S.~Loosemore, R.~M. Stallman, et.\ al., \emph{The GNU
    \code{C}~Library Reference Manual}, GNU Press, 2004.
%
\bibitem{R} K.~Reckdahl, \emph{Using Imported Graphics in \LaTeX2e},
  Version~2.0, whitepaper, Dec.~15, 1997
%
\bibitem{S} B.~Stroustrup, \emph{The \CXX\ Programming Language},
  Special Ed., Addison-Wesley, 1997
%
\bibitem{Z} T. van~Zandt, \emph{PSTricks: PostScript Macros for
    Generic \TeX}, Version~0.93a, whitepaper, Mar.~12, 1993
%
\end{thebibliography}

\clearpage

\printindex

\end{document}

All axis-drawing commands accept an optional second integer argument,
which specifies the number of segments to use when drawing the axis.
This option is useful only when the lens does not map lines to screen
lines.

